% gml.tex - plain TeX macros for the GML backend (xslt/tex.xsl).
%
% Works with Knuth's tex (DVI), pdftex and luatex in plain mode.  No LaTeX.
% The document written by tex.xsl uses only these public commands:
%
%   \gmlbegindocument ... \gmlenddocument      around the whole document
%   \gmlstart{NAME}{k=v,...} ... \gmlend{NAME}  every source element
%       (h1, p, ol, gml:h, gml:fn ...).  \gmlstart opens a TeX group,
%       stores the attribute list (\gmlattr{key}\macro reads it) and calls
%       \gml@NAME@start; \gmlend calls \gml@NAME@end, closes the group and
%       re-applies the formatting state (\gml@apply).  An element without
%       a definition is shown as plain text (message once).
%   \gmlhint{NAME}{k=v,...}       a DCF formatting hint (point event):
%       break=yes calls \gml@break (\par), then \gml@hint@NAME.
%   \gmlbeginlines{WORD}{k=v,...} ... \gmlendlines    a gml:lines scope
%       (.ce .ri .us .up .uc .in/.ir after/for .il .un .rc *); handlers
%       \gml@lines@WORD@pre (before the group) @begin @end @post (after).
%   \gmlnum{label}                gml:num, dispatched on the current
%       element: \gml@num@NAME{label} (h1: label+quad, li: hanging label,
%       figcap: bold hanging label, fn: dropped - the mark is a key).
%   \gmlanchor{name}              a page-reference target
%   \gmlpageref{anchor}{when}{label}   when = always | auto
%   \gmltocbegin{title} \gmltocentry{k=v}{num}{text}{page} \gmltocend
%   \gmlfiglistbegin{title} \gmlfiglistentry{k=v}{label}{text}{page} \gmlfiglistend
%   \gmlindexbegin{title} \gmlindexgroup{letter}
%       \gmlindexentry{level}{term}{refs} \gmlindexend
%       refs: \gmlindexfirst \gmlindexref{anchor}{major=,range=,text=}
%             \gmlindexsep \gmlindexrangesep \gmlsee{kind}{label}{text}
%   \gmlchar{codepoint}           any non-ASCII or TeX-special character
%   \gmltabline{seg\gmltab seg...}  one input line holding tab characters
%   \gmlusword \gmlupword \gmlucword {word}   words inside .us/.up/.uc
%
% Overriding: every internal macro is \gml@...; load your own file after
% this one (a file gmllocal.tex on TEXINPUTS is \input automatically at
% the end of this file), and use the helpers that need no @ catcode:
%   \gmldefstart{h1}{...}  \gmldefend{h1}{...}  \gmldefhint{sp}{...}
%   \gmldeflines{ce}{begin}{...}  \gmldefnum{li}#1{...}
%   \gmldefchar{8364}{EUR}  \gmldeffont{hi1}{\sl}
%   \gmldefhstart{fig}{...}  \gmldefhend{fig}{...}   (hint mode, below)
% or \catcode`\@=11 and redefine \gml@... directly.  A tag set's file
% (\def\gmltagset{file} before \input gml; gmlfmt -tagset writes it) is
% \input at the end too, before gmllocal.tex, so the most local wins.
%
% Hint mode: \def\gmlmode{hints} before \input gml (tex.xsl writes it
% for a document formatted by a real DCF profile, e.g. IBM's DSMPROF3):
% elements are only groups (except gml:h gml:fn gml:rh gml:rf and the fig
% float), \gmlnum prints nothing, and the hints do all the layout.
% DCF definitions: .df (fonts), .dr (named rules, used by .bx and .hr)
% and .dh (head levels) are implemented; see their hints below.
%
% Formatting state.  DCF control words are state (set until changed),
% and so are \leftskip & co. in TeX.  Hints set GLOBAL hint state
% (\gml@hintleft, \gml@hintright, \gml@hintll, fill/align, font, keep);
% elements set LOCAL element state (\gml@elemleft ... and font/fill
% overrides).  \gml@apply combines them: margins add (element indent +
% DCF .in), for fill mode and font the most recent setting wins (a
% global clock stamps each setting), so a .bf inside a paragraph lives
% on after the paragraph and a :hp1 inside a .bf region is still italic.
%
% DCF units (\gml@dimof{n}{unit}\register): ch = width of "0" in the
% body font (\tenrm; cmr10: 5pt); line = \baselineskip; i cm mm; pt = DCF
% (pica) points = 1/72in = TeX bp; dd = Didot points; mh/mv = em of the
% current font (\fontdimen6; DCF's em-height is the font size, so mv is
% also the em, not the ex-height); dh/dv = device units of 1/240in (3820
% pels).  A mixed sum ("4ch+1i", unit=mixed) is parsed term by term.
%
% Page references: \gmlanchor writes \gml@ap{name}{\folio} to
% \jobname.ref at shipout; the file is read by \gmlbegindocument, so
% references are right on the second run ("??" on the first).  A
% when=auto reference also writes its own page (\gml@rp{n}{\folio}) and
% prints " label page" only if the target is on another page (on the
% first run it always prints).  \gmlenddocument re-reads the new file
% and says when TeX must be run again.
%
% Non-ASCII: tex.xsl writes pure ASCII; every character outside
% printable ASCII and every TeX special except $ & # % (which become \$
% \& \# \%) is \gmlchar{decimal code point}.  The table below covers
% Latin-1, Latin Extended-A letters with plain accents, typographic
% punctuation, a few symbols; anything else prints "?" with a message.
% In a typewriter font (\fam=\ttfam) the ASCII specials print as
% themselves.
\ifx\gmlbegindocument\undefined\else\endinput\fi
\chardef\gmlatcode=\catcode`\@
\catcode`\@=11

% ---- registers ---------------------------------------------------------
\newdimen\gml@chwidth \newdimen\gml@devunit
\newdimen\gml@elemleft \newdimen\gml@elemright \newdimen\gml@elemhsize
\newdimen\gml@hintleft \newdimen\gml@hintright \newdimen\gml@hintll
\newdimen\gml@defaulthsize
\newdimen\gml@dima \newdimen\gml@dimb \newdimen\gml@dimc \newdimen\gml@dimd
\newdimen\gml@listind \newdimen\gml@capind
\newdimen\gml@tabpos \newdimen\gml@ofdim
\newdimen\gml@bxleft \newdimen\gml@bxright
\newdimen\gml@basebs
\newdimen\gml@pagelength \newdimen\gml@pagewidth
\newdimen\gml@topmargin \newdimen\gml@botmargin \newdimen\gml@leftmargin
\newcount\gml@clock \newcount\gml@efontstamp \newcount\gml@hfontstamp
\newcount\gml@efillstamp \newcount\gml@hfillstamp
\newcount\gml@cnta \newcount\gml@cntb \newcount\gml@refcount
\newcount\gml@qdepth \newcount\gml@tabidx \newcount\gml@ofstate
\newcount\gmlindexcolumns
\newcount\gml@itemno   % items so far in the current list (local to the list)
\newbox\gml@box \newbox\gml@figbox \newbox\gml@dtbox
\newbox\gml@rhodd \newbox\gml@rheven \newbox\gml@rfodd \newbox\gml@rfeven
\newbox\gml@fnleader \newbox\gml@ixbox
\newwrite\gml@refout \newread\gml@in
\newif\ifgml@plainpage \newif\ifgml@keep \newif\ifgml@aftertitle
\newif\ifgml@userrh \newif\ifgml@userrf \newif\ifgml@rhsup \newif\ifgml@rfsup
\newif\ifgml@rhsuppage \newif\ifgml@rfsuppage
\newif\ifgml@nofolio \newif\ifgml@refschanged \newif\ifgml@reffile
\newif\ifgml@chappage  % the current page begins with an h0/h1 or section heading
\newif\ifgml@tabfound \newif\ifgml@tabfirst \newif\ifgml@rcon
\newif\ifgml@intitlep \newif\ifgml@preserve \newif\ifgml@paoff
\newif\ifgml@found \newif\ifgml@pnpref \newif\ifgml@bxfirst
\newif\ifgml@userleader \newif\ifgml@inhead
\newif\ifgml@hintmode \newif\ifgml@infig \newif\ifgml@fontquiet \newif\ifgml@bxrule
\newdimen\gml@rulewd \newdimen\gml@bxwd \newdimen\gml@hrpos \newdimen\gml@hrstart
\newcount\gml@bxstate \newcount\gml@hrmode

% ---- constants -----------------------------------------------------------
\def\gml@yes{yes} \def\gml@no{no} \def\gml@mixedstr{mixed}
\def\gml@onstr{on} \def\gml@offstr{off} \def\gml@autostr{auto}
\def\gml@inlinestr{inline} \def\gml@boxstr{box} \def\gml@rulestr{rule}
\def\gml@nonestr{none} \def\gml@canstr{can} \def\gml@charstr{char}
\def\gml@cancelstr{cancel} \def\gml@supstr{sup} \def\gml@resstr{res}
\def\gml@resetstr{reset} \def\gml@setstr{set} \def\gml@addstr{add}
\def\gml@delstr{del} \def\gml@defaultstr{default} \def\gml@bystr{by}
\def\gml@normalstr{normal} \def\gml@oddstr{odd} \def\gml@evenstr{even}
\def\gml@leaderstr{leader} \def\gml@topstr{top} \def\gml@offnostr{offno}
\def\gml@romanstr{roman} \def\gml@arabicstr{arabic} \def\gml@prefstr{pref}
\def\gml@nobindstr{nobind} \def\gml@leftstr{left} \def\gml@rightstr{right}
\def\gml@forstr{for} \def\gml@tostr{to} \def\gml@minstr{min} \def\gml@monostr{mono}
\def\gml@sansstr{sans} \def\gml@hintsstr{hints}
\def\gml@nokeymark{\gml@nokey}
\def\gml@end{\gml@end}

% ---- fonts ----------------------------------------------------------------
% Fonts come from a font scheme: a file gmlfont-NAME.tex that fills
% slots with the public commands
%   \gmlfontname{SLOT}{SPEC}  SPEC is what follows \font\x= : a TFM name
%       for tex/pdftex; for luatex/xetex any name the engine (and its
%       font loader) accepts, e.g. "[file.otf]" or "Family:mapping=tex-text".
%       Slots: rm it sl bf bi tt (body, 1x), hd (headings, scaled
%       1.2x-2.488x), and the optional ti (the :title, 2.986x; default:
%       hd) and rm7 rm5 bf7 bf5 for the script sizes (default: rm/bf at
%       7/10 and 5/10 of the body size), ss and ssbf (sans serif and its
%       bold, for .df TYPE('... sans ...'); default: rm and bf).
%   \gmlfontsize{10pt}        body size
% \gmlusefonts{NAME} reads gmlfont-NAME-ENGINE.tex if there is one, else
% gmlfont-NAME.tex (ENGINE = \gmlengine: tex, pdftex, luatex or xetex),
% then (re)loads every font, globally; \gmlusefonts{} only reloads (after
% changing a slot by hand).  Shipped: ncs (New Century Schoolbook, the
% default) and cm (Computer Modern).  Choosing per document or run:
%   gmlfmt -p tex-fonts=NAME   (tex.xsl writes \def\gmlfonts{NAME})
%   \def\gmlfonts{NAME} before \input gml
%   \gmlusefonts{NAME} in gmllocal.tex, tex-preamble, or the source
%       (.xr tex /\gmlusefonts{NAME}/ before the first text)
% Choosing per engine: put a gmlfont-NAME-ENGINE.tex on TEXINPUTS.
% Changing one slot of the scheme (e.g. the title face), in gmllocal.tex,
% tagset.tex or tex-preamble:  \gmlfontname{ti}{pncbi7t}\gmlusefonts{}
\def\gmlengine{tex}
\ifx\pdfoutput\undefined\else\def\gmlengine{pdftex}\fi
\ifx\XeTeXversion\undefined\else\def\gmlengine{xetex}\fi
\ifx\directlua\undefined\else\def\gmlengine{luatex}\fi
\def\gmlfontname#1{\expandafter\def\csname gml@fspec@#1\endcsname}
\def\gmlfontsize#1{\def\gml@fsize{#1}}
\gmlfontsize{10pt}
\def\gmlusefonts#1{\edef\gml@fscheme{#1}%
  \ifx\gml@fscheme\empty\else
    \gml@fclear{rm7}\gml@fclear{rm5}\gml@fclear{bf7}\gml@fclear{bf5}%
    \gml@fclear{ti}\gml@fclear{ss}\gml@fclear{ssbf}%
    \openin\gml@in=gmlfont-\gml@fscheme-\gmlengine.tex\relax
    \ifeof\gml@in\closein\gml@in
      \openin\gml@in=gmlfont-\gml@fscheme.tex\relax
      \ifeof\gml@in\closein\gml@in
        \errmessage{gml.tex: no font scheme gmlfont-\gml@fscheme.tex}%
      \else\closein\gml@in\input gmlfont-\gml@fscheme.tex \fi
    \else\closein\gml@in\input gmlfont-\gml@fscheme-\gmlengine.tex \fi
  \fi\gml@loadfonts}
\def\gml@fclear#1{\expandafter\let\csname gml@fspec@#1\endcsname\relax}
% \gml@fload\cs{slot}{fallback slot}{factor of the body size}
\def\gml@fload#1#2#3#4{\gml@dima=\gml@fsize\relax\gml@dima=#4\gml@dima
  \expandafter\ifx\csname gml@fspec@#2\endcsname\relax
    \edef\gml@fspec{\csname gml@fspec@#3\endcsname}%
  \else\edef\gml@fspec{\csname gml@fspec@#2\endcsname}\fi
  \edef\gml@tmpa{\global\font\noexpand#1=\gml@fspec\space at \the\gml@dima}%
  \gml@tmpa}
\def\gml@loadfonts{%
  \gml@fload\tenrm{rm}{rm}{1}\gml@fload\sevenrm{rm7}{rm}{.7}%
  \gml@fload\fiverm{rm5}{rm}{.5}%
  \gml@fload\tenbf{bf}{bf}{1}\gml@fload\sevenbf{bf7}{bf}{.7}%
  \gml@fload\fivebf{bf5}{bf}{.5}%
  \gml@fload\tenit{it}{it}{1}\gml@fload\tensl{sl}{sl}{1}%
  \gml@fload\tentt{tt}{tt}{1}\gml@fload\gml@tenbi{bi}{bi}{1}%
  \gml@fload\gml@hdfontA{hd}{hd}{2.488}\gml@fload\gml@hdfontB{hd}{hd}{2.074}%
  \gml@fload\gml@hdfontC{hd}{hd}{1.728}\gml@fload\gml@hdfontD{hd}{hd}{1.2}%
  \gml@fload\gml@titlefnt{ti}{hd}{2.986}%
  \gml@fload\gml@tenss{ss}{rm}{1}\gml@fload\gml@tenssbf{ssbf}{bf}{1}%
  \textfont0=\tenrm \scriptfont0=\sevenrm \scriptscriptfont0=\fiverm
  \textfont\bffam=\tenbf \scriptfont\bffam=\sevenbf
  \scriptscriptfont\bffam=\fivebf
  \textfont\itfam=\tenit \textfont\slfam=\tensl \textfont\ttfam=\tentt
  \setbox0\hbox{\tenrm 0}\global\gml@chwidth=\wd0 \rm}
\ifx\gmlfonts\undefined\gmlusefonts{ncs}\else\gmlusefonts{\gmlfonts}\fi
\def\gml@bi{\fam\m@ne\gml@tenbi}
\def\gml@ss{\fam\m@ne\gml@tenss} \def\gml@ssbf{\fam\m@ne\gml@tenssbf}
% \gml@sized{slot}{size}: the slot's face at a DCF point size (TYPE(n) of
% .df; +n/-n relative to 10), scaled with the body: the Starter Set's
% page-printer fonts assume a 10-point body (SH35-0050 p.25), so 9 is
% 9/10 of \gmlfontsize.  The font is loaded on first use, from the
% scheme then in effect.
{\catcode`p=12 \catcode`t=12 \gdef\gml@rmdim#1pt{#1}}
\def\gml@strip#1{\expandafter\gml@rmdim\the#1}
\def\gml@sized#1#2{\gml@sizedim{#2}%
  \expandafter\ifx\csname gml@sf@#1@\number\gml@dima\endcsname\relax
    \expandafter\ifx\csname gml@fspec@#1\endcsname\relax
      \edef\gml@fspec{\csname gml@fspec@\csname gml@fslotfb@#1\endcsname\endcsname}%
    \else\edef\gml@fspec{\csname gml@fspec@#1\endcsname}\fi
    \edef\gml@tmpa{\global\font\expandafter\noexpand
      \csname gml@sf@#1@\number\gml@dima\endcsname=\gml@fspec\space at \number\gml@dima sp}%
    \gml@tmpa\fi
  \fam\m@ne\csname gml@sf@#1@\number\gml@dima\endcsname}
\def\gml@sizedim#1{\gml@sizesign#1\gml@nil
  \gml@dimb=\gml@fsize\relax\divide\gml@dimb by 10
  \gml@dima=\gml@strip\gml@dimc\gml@dimb}
\def\gml@sizesign#1#2\gml@nil{\gml@dimc=#1#2pt\relax
  \if#1+\advance\gml@dimc10pt\fi \if#1-\advance\gml@dimc10pt\fi}
\def\gml@fslotfb@ss{rm} \def\gml@fslotfb@ssbf{bf} \def\gml@fslotfb@ti{hd}
\def\gml@fslotfb@rm7{rm} \def\gml@fslotfb@rm5{rm} \def\gml@fslotfb@bf7{bf}
\def\gml@fslotfb@bf5{bf}
\def\gmldeffont#1{\expandafter\def\csname gml@font@#1\endcsname}
% DCF / Starter Set font names (SH35-0050 Figure 6, page printers)
\gmldeffont{hi0}{\rm}  \gmldeffont{hi1}{\it}
\gmldeffont{hi2}{\bf}  \gmldeffont{hi3}{\gml@bi}
\gmldeffont{hd0}{\fam\m@ne\gml@hdfontA}  \gmldeffont{hd1}{\fam\m@ne\gml@hdfontA}
\gmldeffont{hd2}{\fam\m@ne\gml@hdfontB}  \gmldeffont{hd3}{\fam\m@ne\gml@hdfontC}
\gmldeffont{hd4}{\fam\m@ne\gml@hdfontD}  \gmldeffont{hd5}{\bf}
\gmldeffont{hd6}{\gml@bi}
\gmldeffont{title}{\fam\m@ne\gml@titlefnt}
\gmldeffont{normal}{\rm} \gmldeffont{roman}{\rm} \gmldeffont{bold}{\bf}
\gmldeffont{italic}{\it} \gmldeffont{bolditalic}{\gml@bi}
\gmldeffont{mono}{\tt} \gmldeffont{xmp}{\tt} \gmldeffont{slanted}{\sl}
\gmldeffont{rh}{\rm} \gmldeffont{rf}{\rm} \gmldeffont{figcap}{\rm}
% fonts of DSMPROF3's page-printer section that its line-device sections
% leave undefined (".bf xmpfont =" then restarts the current font): TeX
% is a page device, so the page-printer faces are the defaults here.  A
% .df in the document replaces them.
\gmldeffont{xmpfont}{\tt} \gmldeffont{fnt}{\gml@sized{rm}{9}}
\def\gml@hi#1{\csname gml@font@hi#1\endcsname}

% ---- page layout (DCF page printer defaults, one column) ----------------
\gml@pagelength=11in \gml@pagewidth=8.5in
\gml@topmargin=1in \gml@botmargin=1.1in \gml@leftmargin=1in
\hsize=6.5in \gml@defaulthsize=\hsize \gml@elemhsize=\hsize
\def\gml@layout{%
  \global\vsize\gml@pagelength \global\advance\vsize-\gml@topmargin
  \global\advance\vsize-\gml@botmargin
  \global\voffset\gml@topmargin \global\advance\voffset-1in
  \global\hoffset\gml@leftmargin \global\advance\hoffset-1in
  \ifx\pdfpagewidth\undefined\else
    \global\pdfpagewidth\gml@pagewidth \global\pdfpageheight\gml@pagelength\fi}
\gml@layout
\parindent=\z@ \parskip=\z@
\gml@basebs=12pt \baselineskip=\gml@basebs
\tolerance=1000 \emergencystretch=2em \hbadness=5000 \vbadness=5000
\hyphenpenalty=\@M   % DCF: hyphenation is off until .hy on
\raggedbottom
\gml@devunit=1in \divide\gml@devunit by 240
\gmlindexcolumns=2
\def\gmltitlealign{right}     % SYSVART: right | left | center
\def\gmldocnumlabel{Document Number}
\def\gmlnotelabel{Note:}
\def\gmlabstracttitle{Abstract}
\def\gmlprefacetitle{Preface}
\def\gml@titletop{2in}
\def\gml@titlegapdim{2in}
\def\gml@colsep{2em}
\def\gml@tocpage{2.5em}
\def\gml@figpad{4pt}
\def\gml@rulesep{3pt}
\def\gml@folioprefix{}

% ---- utilities -------------------------------------------------------------
\def\gml@warnonce#1#2{\expandafter\ifx\csname gml@w@#1\endcsname\relax
  \immediate\write16{GML: #2}\expandafter\gdef\csname gml@w@#1\endcsname{}\fi}
\def\gml@notimpl#1{\gml@warnonce{hint#1}{hint .#1 is not implemented in TeX (ignored)}}
% \gml@get\list{key}\macro: \macro := the value of key in list, or empty
\def\gml@get#1#2#3{\def\gml@get@a##1,#2=##2,##3\gml@nil{\def#3{##2}}%
  \expandafter\gml@get@a\expandafter,#1,#2=\gml@nokey,\gml@nil
  \ifx#3\gml@nokeymark\let#3\empty\fi}
\def\gml@aget{\gml@get\gml@attrs}
\def\gml@hget{\gml@get\gml@hattrs}
\def\gml@lget{\gml@get\gml@thislattrs}
\def\gmlattr#1#2{\gml@aget{#1}#2}
\def\gmlhintattr#1#2{\gml@hget{#1}#2}
\let\gml@attrs\empty \let\gml@hattrs\empty \let\gml@lattrs\empty
\let\gml@thislattrs\empty \def\gml@cur{}
% \gml@forwords{list}\cmd: \cmd{word} for each blank-separated word
\def\gml@forwords#1#2{\let\gml@fw@do#2\edef\gml@fw@list{#1\space}%
  \expandafter\gml@fw@a\gml@fw@list\gml@nil}
\def\gml@fw@a#1 #2\gml@nil{\def\gml@fw@w{#1}%
  \ifx\gml@fw@w\empty\else\gml@fw@do{#1}\fi
  \def\gml@fw@w{#2}\ifx\gml@fw@w\empty\else\gml@fw@a#2\gml@nil\fi}
% \gml@ifin{\macro}{string}: \ifgml@found
\def\gml@ifin#1#2{\def\gml@ifin@a##1#2##2\gml@nil{\def\gml@ifin@r{##2}}%
  \expandafter\gml@ifin@a#1#2\gml@nil
  \ifx\gml@ifin@r\empty\gml@foundfalse\else\gml@foundtrue\fi}
\def\gml@lower#1{\edef\gml@tmpl{#1}\lowercase\expandafter{\expandafter
  \def\expandafter\gml@tmpl\expandafter{\gml@tmpl}}}

% ---- DCF units ---------------------------------------------------------------
\def\gml@u@ch#1#2{#2=#1\gml@chwidth}
\def\gml@u@line#1#2{#2=#1\baselineskip}
\def\gml@u@i#1#2{#2=#1in}
\def\gml@u@cm#1#2{#2=#1cm}
\def\gml@u@mm#1#2{#2=#1mm}
\def\gml@u@pt#1#2{#2=#1bp}
\def\gml@u@dd#1#2{#2=#1dd}
\def\gml@u@mh#1#2{#2=#1\fontdimen6\font}
\let\gml@u@mv\gml@u@mh \let\gml@u@m\gml@u@mh
\def\gml@u@dh#1#2{#2=#1\gml@devunit} \let\gml@u@dv\gml@u@dh
\expandafter\let\csname gml@u@\endcsname\gml@u@ch
% \gml@dimof{number}{unit}\register   (unit "mixed": number is a sum)
\def\gml@dimof#1#2#3{\edef\gml@dv{#1}\edef\gml@du{#2}%
  \ifx\gml@du\gml@mixedstr \expandafter\gml@sum\expandafter{\gml@dv}#3%
  \else\expandafter\ifx\csname gml@u@\gml@du\endcsname\relax
    \gml@warnonce{unit\gml@du}{unknown space unit \gml@du\space(taken as ch)}%
    \gml@u@ch\gml@dv#3%
  \else\csname gml@u@\gml@du\endcsname\gml@dv#3\fi\fi}
% \gml@sum{4ch+1i}\register; also one token "15pt" (\gml@dimtok)
\def\gml@sum#1#2{\let\gml@sumreg#2\gml@sumreg=\z@
  \def\gml@sn{}\def\gml@su{}\gml@sum@a#1\gml@end}
\def\gml@sum@a#1{%
  \ifx#1\gml@end \gml@sum@flush \let\gml@next\relax
  \else \let\gml@next\gml@sum@a
    \ifx#1+\gml@sum@flush\def\gml@sn{}%
    \else\ifx#1-\gml@sum@flush\def\gml@sn{-}%
    \else\ifnum`#1<`a\relax
      \ifx\gml@su\empty\else\gml@sum@flush\fi\edef\gml@sn{\gml@sn#1}%
    \else\edef\gml@su{\gml@su#1}\fi\fi\fi
  \fi\gml@next}
\def\gml@sum@flush{\ifx\gml@sn\empty\else
  \ifx\gml@su\empty\def\gml@su{ch}\fi
  \gml@dimof\gml@sn\gml@su\gml@dimd \advance\gml@sumreg\gml@dimd
  \def\gml@sn{}\def\gml@su{}\fi}
\def\gml@dimtok#1#2{\edef\gml@tmpt{#1}\expandafter\gml@sum\expandafter{\gml@tmpt}#2}
% the number part of a token ("2line" -> 2) in \gml@sn
\def\gml@numpart#1{\def\gml@sn{}\edef\gml@tmpt{#1}\expandafter\gml@np@a\gml@tmpt\gml@end}
\def\gml@np@a#1{\ifx#1\gml@end\let\gml@next\relax\else\let\gml@next\gml@np@a
  \ifnum`#1<`a\relax\edef\gml@sn{\gml@sn#1}\fi\fi\gml@next}
% a hint dimension: \gml@hdim{attr}{unit attr}\register
\def\gml@hdim#1#2#3{\gml@hget{#1}\gml@tv\gml@hget{#2}\gml@tu\gml@dimof\gml@tv\gml@tu#3}

% ---- formatting state ----------------------------------------------------------
\def\gml@efont{\rm} \let\gml@hfont\empty
\let\gml@efill\empty \let\gml@ealign\empty
\def\gml@hfill{yes} \def\gml@halign{justify}
\def\gml@setfont#1{\global\advance\gml@clock\@ne \gml@efontstamp\gml@clock
  \def\gml@efont{#1}#1}
\def\gml@setfill#1#2{\global\advance\gml@clock\@ne \gml@efillstamp\gml@clock
  \edef\gml@efill{#1}\edef\gml@ealign{#2}}
\def\gml@hsetfill#1#2{\global\advance\gml@clock\@ne \global\gml@hfillstamp\gml@clock
  \xdef\gml@hfill{#1}\xdef\gml@halign{#2}}
\def\gml@hsetfontcmd{\global\advance\gml@clock\@ne \global\gml@hfontstamp\gml@clock
  \global\let\gml@hfont\gml@fontcmd}
\def\gml@applyfont{\ifx\gml@hfont\empty\gml@efont
  \else\ifnum\gml@efontstamp>\gml@hfontstamp\gml@efont\else\gml@hfont\fi\fi}
\def\gml@applyfill{%
  \ifx\gml@efill\empty \let\gml@curfill\gml@hfill \let\gml@curalign\gml@halign
  \else\ifnum\gml@efillstamp>\gml@hfillstamp
    \let\gml@curfill\gml@efill \let\gml@curalign\gml@ealign
  \else \let\gml@curfill\gml@hfill \let\gml@curalign\gml@halign\fi\fi
  \ifx\gml@curfill\gml@no\catcode`\^^M=\active\else\catcode`\^^M=5 \fi}
\def\gml@applymargins{%
  \gml@dima\gml@elemleft \advance\gml@dima\gml@hintleft
  \gml@dimb\gml@elemright \advance\gml@dimb\gml@hintright
  \leftskip\gml@dima \rightskip\gml@dimb \parfillskip\z@ plus1fil
  \csname gml@align@\gml@curalign\endcsname}
\def\gml@align@justify{\ifx\gml@curfill\gml@no\advance\rightskip\z@ plus1fil\fi}
\def\gml@align@left{\ifx\gml@curfill\gml@no\advance\rightskip\z@ plus1fil
  \else\advance\rightskip\z@ plus3em\fi}
\def\gml@align@right{\advance\leftskip\z@ plus1fil \parfillskip\z@}
\def\gml@align@center{\advance\leftskip\z@ plus1fil
  \advance\rightskip\z@ plus1fil \parfillskip\z@}
\let\gml@align@inside\gml@align@left   % no page parity at typesetting time
\let\gml@align@outside\gml@align@right
\def\gml@applyhsize{\ifdim\gml@hintll>\z@\hsize\gml@hintll\else\hsize\gml@elemhsize\fi}
\def\gml@applykeep{\ifgml@keep\interlinepenalty\@M\else\interlinepenalty\z@\fi}
\def\gml@apply{\gml@applyfont\gml@applyfill\gml@applymargins
  \gml@applyhsize\gml@applykeep}
% an obeyed end of line (format off): the line is a paragraph; an empty
% line is a blank line only in xml:space="preserve" text
{\catcode`\^^M=\active \gdef^^M{\gml@obeyedline}}
\def\gml@obeyedline{\ifhmode\par\else\ifgml@preserve\leavevmode\par\fi\fi}
% DCF .of and revision codes act on every paragraph start
\def\gml@everypar{\ifcase\gml@ofstate\or
    \hangindent\gml@ofdim\hangafter\@ne\global\gml@ofstate\tw@
  \or\hangindent\gml@ofdim\hangafter\z@\fi
  \ifgml@rcon\llap{\gml@rcchar\enspace}\fi}
\def\gml@rcchar{|}

% ---- vertical space ------------------------------------------------------------
% \gml@vskip: a DCF skip (.sk): discarded at the top of a page, and
% consecutive skips do not add up (the larger one counts).
\def\gml@vskip#1{\par\ifvmode\gml@dimd=#1\relax
  \ifdim\lastskip<\gml@dimd \ifdim\lastskip>\z@\vskip-\lastskip\fi
    \ifgml@keep\nobreak\fi\vskip\gml@dimd\fi\fi}
% \gml@vspace: a DCF space (.sp): kept at the top of a page
\def\gml@vspace#1{\par\ifvmode\vglue#1\relax\fi}
\def\gml@newpage{\par\vfill\penalty-\@M}
% \gml@chapnewpage: a new page for an h0/h1 or section heading.  Pending
% floats are flushed first (\supereject: plain's \dosupereject puts them on
% pages of their own), so a figure of the previous chapter never lands above
% the new heading; the page is then marked as a chapter page (see fig).
% The page builder runs synchronously here, so \output has already reset
% the flag for the pages shipped before it is set.
\def\gml@chapnewpage{\par\vfill\supereject\global\gml@chappagetrue}
\def\gml@needspace#1{\par\ifvmode\ifinner\else\gml@dimc=#1\relax\penalty\z@
  \ifdim\pagegoal<\maxdimen \gml@dimb\pagegoal\advance\gml@dimb-\pagetotal
    \ifdim\gml@dimb<\gml@dimc\vfill\penalty-\@M\fi\fi\fi\fi}
\def\gml@parstart#1{\ifhmode\ifdim\lastkern=1sp\unkern\else\gml@vskip{#1}\fi
  \else\gml@vskip{#1}\fi}

% ---- dispatch --------------------------------------------------------------------
% Hint mode (\def\gmlmode{hints} before \input gml; tex.xsl writes it for
% a document formatted by a real DCF profile): the profile's control
% words (hints) do the layout, so an element calls \gml@h@NAME@start /
% \gml@h@NAME@end if defined (\gmldefhstart, \gmldefhend) and is
% otherwise only a group; \gmlnum prints nothing (the profile writes its
% own numbers).  Defined: gml:h gml:fn gml:rh gml:rf (as in element
% mode) and fig (the float and its frame, without margins).
\def\gmlstart#1#2{\begingroup\def\gml@cur{#1}\def\gml@attrs{#2}%
  \ifgml@hintmode\gml@hcall{#1}{start}%
  \else\expandafter\ifx\csname gml@#1@start\endcsname\relax
    \gml@warnonce{elt#1}{no TeX style for element #1 (shown as plain text)}%
  \else\csname gml@#1@start\endcsname\fi\fi}
\def\gmlend#1{\ifgml@hintmode\gml@hcall{#1}{end}%
  \else\expandafter\ifx\csname gml@#1@end\endcsname\relax
  \else\csname gml@#1@end\endcsname\fi\fi\endgroup\gml@apply}
\def\gml@hcall#1#2{\expandafter\ifx\csname gml@h@#1@#2\endcsname\relax
  \else\csname gml@h@#1@#2\endcsname\fi}
\def\gmldefhstart#1{\expandafter\def\csname gml@h@#1@start\endcsname}
\def\gmldefhend#1{\expandafter\def\csname gml@h@#1@end\endcsname}
% hint mode: these elements keep their element-mode style
\def\gml@hkeep#1{\gmldefhstart{#1}{\csname gml@#1@start\endcsname}%
  \gmldefhend{#1}{\csname gml@#1@end\endcsname}}
\gml@hkeep{gml:h} \gml@hkeep{gml:fn} \gml@hkeep{gml:rh} \gml@hkeep{gml:rf}
\ifx\gmlmode\undefined\else\ifx\gmlmode\gml@hintsstr\gml@hintmodetrue\fi\fi
\def\gmldefstart#1{\expandafter\def\csname gml@#1@start\endcsname}
\def\gmldefend#1{\expandafter\def\csname gml@#1@end\endcsname}
\def\gml@letelement#1#2{%
  \expandafter\let\csname gml@#1@start\expandafter\endcsname\csname gml@#2@start\endcsname
  \expandafter\let\csname gml@#1@end\expandafter\endcsname\csname gml@#2@end\endcsname}
\def\gmldefhint#1{\expandafter\def\csname gml@hint@#1\endcsname}
\def\gmldeflines#1#2{\expandafter\def\csname gml@lines@#1@#2\endcsname}
\def\gmldefnum#1{\expandafter\def\csname gml@num@#1\endcsname}
\def\gmldefchar#1{\expandafter\def\csname gml@char@#1\endcsname}

\def\gml@break{\par}
\def\gmlhint#1#2{\def\gml@hattrs{#2}\gml@hget{break}\gml@tmpb
  \ifx\gml@tmpb\gml@yes\gml@break\fi
  \expandafter\ifx\csname gml@hint@#1\endcsname\relax\gml@notimpl{#1}%
  \else\csname gml@hint@#1\endcsname\fi}

\def\gmlbeginlines#1#2{\def\gml@lattrs{#2}\gml@get\gml@lattrs{break}\gml@tmpb
  \ifx\gml@tmpb\gml@yes\gml@break\fi
  \expandafter\ifx\csname gml@lines@#1@begin\endcsname\relax
    \gml@warnonce{lines#1}{line scope .#1 is not implemented in TeX (ignored)}\fi
  \gml@lcall{#1}{pre}\begingroup\def\gml@lword{#1}\let\gml@thislattrs\gml@lattrs
  \gml@lcall{#1}{begin}}
\def\gml@lcall#1#2{\expandafter\ifx\csname gml@lines@#1@#2\endcsname\relax
  \else\csname gml@lines@#1@#2\endcsname\fi}
\def\gmlendlines{\gml@lcall\gml@lword{end}%
  \expandafter\gml@endlines@a\expandafter{\gml@lword}}
\def\gml@endlines@a#1{\endgroup\gml@apply\gml@lcall{#1}{post}}
\def\gml@lmore{\gml@lget{more}\gml@tmpb\ifx\gml@tmpb\gml@yes\else\ifhmode\par\fi\fi}

\def\gmlnum#1{\ifgml@hintmode\else
  \expandafter\ifx\csname gml@num@\gml@cur\endcsname\relax
  #1\enspace\else\csname gml@num@\gml@cur\endcsname{#1}\fi\fi}

% ---- characters ---------------------------------------------------------------------
\def\gml@char#1{\expandafter\ifx\csname gml@char@#1\endcsname\relax
  \gml@warnonce{char#1}{no rendering for character #1 (decimal code point), printed as ?}?%
  \else\csname gml@char@#1\endcsname\fi}
\let\gmlchar\gml@char
\def\gml@ttor#1#2{\ifnum\fam=\ttfam\char#1\relax\else#2\fi}
\gmldefchar{34}{\gml@ttor{34}{''}}
\gmldefchar{60}{\gml@ttor{60}{$<$}}
\gmldefchar{62}{\gml@ttor{62}{$>$}}
\gmldefchar{92}{\gml@ttor{92}{$\backslash$}}
\gmldefchar{94}{\gml@ttor{94}{\^{}}}
\gmldefchar{95}{\gml@ttor{95}{\_}}
\gmldefchar{123}{\gml@ttor{123}{$\{$}}
\gmldefchar{124}{\gml@ttor{124}{$\vert$}}
\gmldefchar{125}{\gml@ttor{125}{$\}$}}
\gmldefchar{126}{\gml@ttor{126}{\~{}}}
\gmldefchar{44}{,} \gmldefchar{61}{=}
% Latin-1
\gmldefchar{160}{\nobreak\ }        \gmldefchar{161}{!`}
\gmldefchar{162}{\rlap/c}            \gmldefchar{163}{{\it\$}}
\gmldefchar{164}{$\diamond$}         \gmldefchar{165}{Y\llap=}
\gmldefchar{166}{$\vert$}            \gmldefchar{167}{\S}
\gmldefchar{168}{\"{}}               \gmldefchar{169}{\copyright}
\gmldefchar{170}{$^{\underline{\rm a}}$} \gmldefchar{171}{$\ll$}
\gmldefchar{172}{$\neg$}             \gmldefchar{173}{\-}
\gmldefchar{174}{{\ooalign{\hfil\raise.07ex\hbox{\sevenrm R}\hfil\crcr\mathhexbox20D}}}
\gmldefchar{175}{\={}}               \gmldefchar{176}{$^\circ$}
\gmldefchar{177}{$\pm$}              \gmldefchar{178}{$^2$}
\gmldefchar{179}{$^3$}               \gmldefchar{180}{\'{}}
\gmldefchar{181}{$\mu$}              \gmldefchar{182}{\P}
\gmldefchar{183}{$\cdot$}            \gmldefchar{184}{\c{}}
\gmldefchar{185}{$^1$}               \gmldefchar{186}{$^{\underline{\rm o}}$}
\gmldefchar{187}{$\gg$}              \gmldefchar{188}{$1\over4$}
\gmldefchar{189}{$1\over2$}          \gmldefchar{190}{$3\over4$}
\gmldefchar{191}{?`}
\gmldefchar{192}{\`A} \gmldefchar{193}{\'A} \gmldefchar{194}{\^A} \gmldefchar{195}{\~A}
\gmldefchar{196}{\"A} \gmldefchar{197}{\AA} \gmldefchar{198}{\AE} \gmldefchar{199}{\c C}
\gmldefchar{200}{\`E} \gmldefchar{201}{\'E} \gmldefchar{202}{\^E} \gmldefchar{203}{\"E}
\gmldefchar{204}{\`I} \gmldefchar{205}{\'I} \gmldefchar{206}{\^I} \gmldefchar{207}{\"I}
\gmldefchar{208}{D\llap{\raise.25ex\hbox{\vrule height.4pt width.35em}\kern.35em}}
\gmldefchar{209}{\~N} \gmldefchar{210}{\`O} \gmldefchar{211}{\'O} \gmldefchar{212}{\^O}
\gmldefchar{213}{\~O} \gmldefchar{214}{\"O} \gmldefchar{215}{$\times$} \gmldefchar{216}{\O}
\gmldefchar{217}{\`U} \gmldefchar{218}{\'U} \gmldefchar{219}{\^U} \gmldefchar{220}{\"U}
\gmldefchar{221}{\'Y} \gmldefchar{222}{Th} \gmldefchar{223}{\ss}
\gmldefchar{224}{\`a} \gmldefchar{225}{\'a} \gmldefchar{226}{\^a} \gmldefchar{227}{\~a}
\gmldefchar{228}{\"a} \gmldefchar{229}{\aa} \gmldefchar{230}{\ae} \gmldefchar{231}{\c c}
\gmldefchar{232}{\`e} \gmldefchar{233}{\'e} \gmldefchar{234}{\^e} \gmldefchar{235}{\"e}
\gmldefchar{236}{\`\i} \gmldefchar{237}{\'\i} \gmldefchar{238}{\^\i} \gmldefchar{239}{\"\i}
\gmldefchar{240}{dh} \gmldefchar{241}{\~n} \gmldefchar{242}{\`o} \gmldefchar{243}{\'o}
\gmldefchar{244}{\^o} \gmldefchar{245}{\~o} \gmldefchar{246}{\"o} \gmldefchar{247}{$\div$}
\gmldefchar{248}{\o} \gmldefchar{249}{\`u} \gmldefchar{250}{\'u} \gmldefchar{251}{\^u}
\gmldefchar{252}{\"u} \gmldefchar{253}{\'y} \gmldefchar{254}{th} \gmldefchar{255}{\"y}
% Latin Extended-A (a selection)
\gmldefchar{262}{\'C} \gmldefchar{263}{\'c} \gmldefchar{268}{\v C} \gmldefchar{269}{\v c}
\gmldefchar{282}{\v E} \gmldefchar{283}{\v e} \gmldefchar{286}{\u G} \gmldefchar{287}{\u g}
\gmldefchar{304}{\.I} \gmldefchar{305}{\i} \gmldefchar{321}{\L} \gmldefchar{322}{\l}
\gmldefchar{323}{\'N} \gmldefchar{324}{\'n} \gmldefchar{327}{\v N} \gmldefchar{328}{\v n}
\gmldefchar{336}{\H O} \gmldefchar{337}{\H o} \gmldefchar{338}{\OE} \gmldefchar{339}{\oe}
\gmldefchar{344}{\v R} \gmldefchar{345}{\v r} \gmldefchar{346}{\'S} \gmldefchar{347}{\'s}
\gmldefchar{350}{\c S} \gmldefchar{351}{\c s} \gmldefchar{352}{\v S} \gmldefchar{353}{\v s}
\gmldefchar{366}{\r U} \gmldefchar{367}{\r u} \gmldefchar{368}{\H U} \gmldefchar{369}{\H u}
\gmldefchar{376}{\"Y} \gmldefchar{377}{\'Z} \gmldefchar{378}{\'z} \gmldefchar{379}{\.Z}
\gmldefchar{380}{\.z} \gmldefchar{381}{\v Z} \gmldefchar{382}{\v z}
\gmldefchar{710}{\^{}} \gmldefchar{732}{\~{}}
% punctuation and symbols
\gmldefchar{8194}{\enspace} \gmldefchar{8195}{\quad} \gmldefchar{8201}{\thinspace}
\gmldefchar{8203}{} \gmldefchar{8208}{-} \gmldefchar{8209}{\hbox{-}}
\gmldefchar{8211}{--} \gmldefchar{8212}{---} \gmldefchar{8213}{---}
\gmldefchar{8216}{`} \gmldefchar{8217}{'} \gmldefchar{8218}{,}
\gmldefchar{8220}{``} \gmldefchar{8221}{''} \gmldefchar{8222}{,,}
\gmldefchar{8224}{$\dagger$} \gmldefchar{8225}{$\ddagger$} \gmldefchar{8226}{$\bullet$}
\gmldefchar{8230}{\dots} \gmldefchar{8240}{\%\raise.2ex\hbox{\sevenrm0}}
\gmldefchar{8242}{$'$} \gmldefchar{8243}{$''$}
\gmldefchar{8249}{$\langle$} \gmldefchar{8250}{$\rangle$}
\gmldefchar{8364}{EUR} \gmldefchar{8482}{$^{\rm TM}$}
\gmldefchar{8592}{$\leftarrow$} \gmldefchar{8594}{$\rightarrow$}
\gmldefchar{8593}{$\uparrow$} \gmldefchar{8595}{$\downarrow$}
\gmldefchar{8656}{$\Leftarrow$} \gmldefchar{8658}{$\Rightarrow$}
\gmldefchar{8722}{$-$} \gmldefchar{8730}{$\surd$} \gmldefchar{8734}{$\infty$}
\gmldefchar{8800}{$\ne$} \gmldefchar{8804}{$\le$} \gmldefchar{8805}{$\ge$}
\gmldefchar{8776}{$\approx$} \gmldefchar{9632}{$\blacksquare$}
\gmldefchar{9633}{$\square$} \gmldefchar{9675}{$\circ$} \gmldefchar{9679}{$\bullet$}
% upper case (.up/.uc scopes)
\def\gml@ucchar#1{\gml@cnta=#1\relax
  \ifnum\gml@cnta>223 \ifnum\gml@cnta<255 \ifnum\gml@cnta=247 \else
    \advance\gml@cnta-32 \fi\fi\fi
  \ifnum\gml@cnta=255 \gml@cnta=376 \fi
  \ifnum\gml@cnta>255 \ifnum\gml@cnta<383 \ifodd\gml@cnta
    \ifnum\gml@cnta=305 \else\ifnum\gml@cnta<312 \advance\gml@cnta\m@ne
    \else\ifnum\gml@cnta>328 \advance\gml@cnta\m@ne\fi\fi\fi
  \else\ifnum\gml@cnta>312 \ifnum\gml@cnta<329 \advance\gml@cnta\m@ne\fi
    \else\ifnum\gml@cnta>377 \advance\gml@cnta\m@ne\fi\fi\fi\fi\fi
  \expandafter\gml@char\expandafter{\the\gml@cnta}}
\def\gmlupword#1{{\let\gmlchar\gml@ucchar\uppercase{#1}}}
\def\gmlusword#1{\underbar{#1}}
\def\gmlucword#1{\underbar{\gmlupword{#1}}}

% ---- tabs (.tp/.tb) -----------------------------------------------------------------
% A line with tab characters is \gmltabline{seg\gmltab seg...}.  The
% segments are boxed and measured from the start of the line; the k-th
% tab goes to the first stop after the previous one that lies beyond
% the current position (DCF order: the stops are used as listed).
\let\gml@tabstops\empty
\def\gml@tabadd{\gml@forwords\gml@tu\gml@tab@addone}
\def\gml@tab@addone#1{\gml@tabsplit#1//\gml@nil{#1}}
\def\gml@tabsplit#1/#2/#3\gml@nil#4{\def\gml@tmpd{#2}%
  \ifx\gml@tmpd\empty\gml@dimtok{#1}\gml@dima\toks2{}%
  \else\gml@dimtok{#2}\gml@dima\toks2{#1}\fi
  \toks@\expandafter{\gml@tabstops}\toks4{#4}%
  \xdef\gml@tabstops{\the\toks@\noexpand\gml@ts{\number\gml@dima}{\the\toks2}{\the\toks4}}}
\def\gml@tabdel{\gml@forwords\gml@tu\gml@tab@delone}
\def\gml@tab@delone#1{\def\gml@deltok{#1}\let\gml@newstops\empty
  \let\gml@ts\gml@ts@keep\gml@tabstops\global\let\gml@tabstops\gml@newstops}
\def\gml@ts@keep#1#2#3{\def\gml@tmpd{#3}\ifx\gml@tmpd\gml@deltok\else
  \toks@\expandafter{\gml@newstops}\toks2{#2}\toks4{#3}%
  \edef\gml@newstops{\the\toks@\noexpand\gml@ts{#1}{\the\toks2}{\the\toks4}}\fi}
\def\gml@tabdefault{\global\let\gml@tabstops\empty
  \def\gml@tu{5ch 10ch 15ch 20ch 25ch 30ch 35ch 40ch 45ch 50ch 55ch 60ch
    65ch 70ch 75ch 80ch}\gml@tabadd}
\def\gml@ts@find#1#2#3{\advance\gml@cnta\@ne
  \ifgml@tabfound\else\ifnum\gml@cnta>\gml@tabidx \ifdim#1sp>\gml@tabpos
    \gml@tabfoundtrue\gml@tabidx\gml@cnta\gml@dimb#1sp\def\gml@tabfill{#2}\fi\fi\fi}
\def\gml@tabadvance{\gml@cnta\z@\gml@tabfoundfalse\let\gml@ts\gml@ts@find\gml@tabstops
  \ifgml@tabfound \advance\gml@dimb-\gml@tabpos
    \ifx\gml@tabfill\empty\hskip\gml@dimb
    \else\hbox to\gml@dimb{\xleaders\hbox{\gml@tabfill}\hfil}\fi
    \advance\gml@tabpos\gml@dimb
  \else\ \fi}
\def\gmltab{\gml@warnonce{tab}{tab outside \string\gmltabline}\ }
\def\gmltabline#1{\leavevmode\begingroup\gml@tabpos\z@\gml@tabidx\z@\gml@tabfirsttrue
  \gml@tabloop#1\gmltab\gml@nil\endgroup}
\def\gml@tabloop#1\gmltab#2\gml@nil{\setbox\gml@box\hbox{#1}%
  \ifgml@tabfirst\gml@tabfirstfalse\else\gml@tabadvance\fi
  \advance\gml@tabpos\wd\gml@box\box\gml@box
  \def\gml@tmpd{#2}\ifx\gml@tmpd\empty\else\gml@tabloop#2\gml@nil\fi}

% ---- parshape for temporary indents ---------------------------------------------------
% \gml@shape{a}{n}{dl}{dr}: a lines unchanged, n lines moved by dl/dr
% (relative to the current margins), then unchanged
\def\gml@shape#1#2#3#4{\gml@dima=#3\relax\gml@dimb=#4\relax
  \gml@dimc\hsize\advance\gml@dimc-\gml@dima\advance\gml@dimc-\gml@dimb
  \let\gml@shapelist\empty
  \count@=#1\relax\loop\ifnum\count@>\z@
    \edef\gml@shapelist{\gml@shapelist\space0pt\space\the\hsize}\advance\count@\m@ne\repeat
  \count@=#2\relax\loop\ifnum\count@>\z@
    \edef\gml@shapelist{\gml@shapelist\space\the\gml@dima\space\the\gml@dimc}\advance\count@\m@ne\repeat
  \gml@cnta=#1\relax\advance\gml@cnta#2\relax\advance\gml@cnta\@ne
  \parshape\gml@cnta\gml@shapelist\space0pt\space\hsize\relax}
% lines in a vertical amount: \gml@cntb := ceiling(dim / baselineskip)
\def\gml@linesof#1{\gml@dimd=#1\relax\gml@cntb\gml@dimd
  \gml@cnta\baselineskip\divide\gml@cntb\gml@cnta
  \multiply\gml@cnta\gml@cntb\ifnum\gml@cnta<\gml@dimd\advance\gml@cntb\@ne\fi}

% ---- hints ------------------------------------------------------------------------
\gmldefhint{br}{}
\gmldefhint{sp}{\gml@hdim{n}{unit}\gml@dima\gml@vspace\gml@dima}
\gmldefhint{sk}{\gml@hdim{n}{unit}\gml@dima\gml@vskip\gml@dima}
% .in/.ir: permanent margins; temporary (FOR/AFTER not in lines): parshape
% a permanent .in that ends the offset carries offset=0 (SH35-0069 p.80)
\def\gml@indent#1#2{\gml@hget{temporary}\gml@tmpb
  \ifx\gml@tmpb\gml@yes\gml@tempindent{#1}{#2}%
  \else\gml@hdim{#1}{unit}\gml@dima\global#2\gml@dima
    \gml@hget{offset}\gml@tmpb
    \ifx\gml@tmpb\empty\else\global\gml@ofdim\z@\global\gml@ofstate\z@\fi
    \gml@apply\fi}
\gmldefhint{in}{\gml@indent{left}\gml@hintleft}
\gmldefhint{ir}{\gml@indent{right}\gml@hintright}
\def\gml@tempindent#1#2{\gml@hdim{#1}{unit}\gml@dima\advance\gml@dima-#2%
  \edef\gml@tdelta{\the\gml@dima}%
  \gml@hget{for}\gml@tv\gml@hget{for-unit}\gml@tu
  \ifx\gml@tv\empty\def\gml@tfor{1}\else
    \gml@dimof\gml@tv\gml@tu\gml@dimc\gml@linesof\gml@dimc\edef\gml@tfor{\the\gml@cntb}\fi
  \gml@hget{after}\gml@tv\gml@hget{after-unit}\gml@tu
  \ifx\gml@tv\empty\def\gml@tafter{0}\else
    \gml@dimof\gml@tv\gml@tu\gml@dimc\gml@linesof\gml@dimc\edef\gml@tafter{\the\gml@cntb}\fi
  \def\gml@tmpb{left}\edef\gml@tmpc{#1}%
  \ifx\gml@tmpb\gml@tmpc\gml@shape\gml@tafter\gml@tfor\gml@tdelta\z@
  \else\gml@shape\gml@tafter\gml@tfor\z@\gml@tdelta\fi}
\gmldefhint{of}{\gml@hdim{offset}{unit}\gml@dima\global\gml@ofdim\gml@dima
  \ifdim\gml@ofdim=\z@\global\gml@ofstate\z@\else\global\gml@ofstate\@ne\fi}
\gmldefhint{ll}{\gml@hget{length}\gml@tv
  \ifx\gml@tv\gml@defaultstr \global\gml@hintll\z@
    \gml@hget{delta}\gml@tv\ifx\gml@tv\empty\else
      \gml@dimtok\gml@tv\gml@dima\global\gml@hintll\gml@defaulthsize
      \global\advance\gml@hintll\gml@dima\fi
  \else\gml@hdim{length}{unit}\gml@dima\global\gml@hintll\gml@dima\fi\gml@apply}
\gmldefhint{cl}{\gml@hget{spec}\gml@tv\ifx\gml@tv\empty\global\gml@hintll\z@
  \else\gml@dimtok\gml@tv\gml@dima\global\gml@hintll\gml@dima\fi\gml@apply}
% format off keeps blank input lines (tex.xsl writes the text as preserved)
\gmldefhint{fo}{\gml@hget{fill}\gml@tv\gml@hget{align}\gml@tu
  \ifx\gml@tv\gml@no\global\gml@preservetrue
  \else\ifx\gml@tv\gml@yes\global\gml@preservefalse\fi\fi
  \gml@hsetfill\gml@tv\gml@tu\gml@apply}
\let\gml@hint@nf\gml@hint@fo \let\gml@hint@ju\gml@hint@fo
\def\gml@blankpage{\hbox{}\vfill\penalty-\@M}
\gmldefhint{pa}{\gml@hget{mode}\gml@tv
  \ifx\gml@tv\gml@offstr\global\gml@paofftrue
  \else\ifx\gml@tv\gml@onstr\global\gml@paofffalse
  \else\ifgml@paoff\else
    \gml@hget{nobreak}\gml@tmpb
    \ifx\gml@tmpb\gml@yes\ifhmode\vadjust{\vfill\penalty-\@M}\else\gml@newpage\fi
    \else\gml@newpage\fi
    \gml@hget{side}\gml@tv
    \ifx\gml@tv\gml@oddstr\ifodd\pageno\else\gml@blankpage\fi\fi
    \ifx\gml@tv\gml@evenstr\ifodd\pageno\gml@blankpage\fi\fi
    \gml@hget{page}\gml@tv\ifx\gml@tv\empty\else\global\pageno=\gml@tv\relax\fi
  \fi\fi\fi}
\gmldefhint{cp}{\gml@hget{n}\gml@tv\ifx\gml@tv\empty\gml@newpage
  \else\gml@hdim{n}{unit}\gml@dima\gml@needspace\gml@dima\fi}
\let\gml@hint@cc\gml@hint@cp
\gmldefhint{cb}{\gml@newpage}
% ---- named rules (.dr) ----------------------------------------------------
% \gml@rule@NAME holds the weight.  boxrule, the default of .bx and .hr,
% is .3mm on page printers (SH35-0069 p.159).  .dr NAME WEIGHT h defines
% one (tex.xsl passes rule= and weight=); DCF-VERIFY: a .dr without a
% weight (DSMPROF3's ".dr @figrule") is taken to have boxrule's weight.
% The FONT operand is for line devices and ignored.
\expandafter\def\csname gml@rule@boxrule\endcsname{.3mm}
\gmldefhint{dr}{\gml@hget{rule}\gml@tv\gml@lower\gml@tv\gml@hget{weight}\gml@tu
  \ifx\gml@tmpl\empty\else
    \ifx\gml@tu\empty\edef\gml@tu{\csname gml@rule@boxrule\endcsname}\fi
    \gml@dimtok\gml@tu\gml@dima
    \expandafter\xdef\csname gml@rule@\gml@tmpl\endcsname{\the\gml@dima}\fi}
% \gml@ruleweight{name}: \gml@rulewd := its weight (boxrule if unknown)
\def\gml@ruleweight#1{\gml@lower{#1}%
  \expandafter\ifx\csname gml@rule@\gml@tmpl\endcsname\relax
    \gml@warnonce{rule#1}{unknown rule #1 (boxrule used)}%
    \gml@rulewd=\csname gml@rule@boxrule\endcsname\relax
  \else\gml@rulewd=\csname gml@rule@\gml@tmpl\endcsname\relax\fi}
% \gml@isnumtok{word}: \ifgml@found if it begins like a number (0-9 + - .)
\def\gml@isnumtok#1{\gml@isnum@a#1\gml@nil}
\def\gml@isnum@a#1#2\gml@nil{\gml@foundfalse
  \expandafter\ifcat\noexpand#1\relax\else
    \ifnum`#1>47 \ifnum`#1<58 \gml@foundtrue\fi\fi
    \if#1+\gml@foundtrue\fi\if#1-\gml@foundtrue\fi\if#1.\gml@foundtrue\fi\fi}
% \gml@kwtest{list}{word}: \ifgml@found if word is one of ,list,
\def\gml@kwtest#1#2{\edef\gml@kwl{,#1,}\gml@ifin\gml@kwl{,#2,}}
% .bx [NEW|SET|ADD|OFF|CAN|CHAR c] [rulename] [LEFT|h|RIGHT|/]... :
% horizontal rules only, from the first to the last position (relative to
% the current left margin; LEFT/RIGHT are the column edges), in the
% weight of the first rule name; vertical rules are not drawn.  .bx
% alone draws a rule across an open box, .bx off closes it (a bottom
% rule), SET opens it without the top rule.  Inside a fig the fig element
% draws its own frame, so the .bx words there (and the .bx off that ends
% that box, after the fig) draw nothing.
\def\gml@bx@tok#1{\gml@isnumtok{#1}%
  \ifgml@found\gml@dimtok{#1}\gml@dima\gml@bx@pos
  \else\def\gml@tmpd{#1}%
    \ifx\gml@tmpd\gml@leftstr\gml@dima-\leftskip\gml@bx@pos
    \else\ifx\gml@tmpd\gml@rightstr\gml@dima\hsize\advance\gml@dima-\leftskip\gml@bx@pos
    \else\gml@kwtest{new,set,add,on,off,can,char,/}{#1}%
      \ifgml@found\else\ifgml@bxrule\else
        \gml@bxruletrue\gml@ruleweight{#1}\global\gml@bxwd\gml@rulewd\fi\fi
    \fi\fi\fi}
\def\gml@bx@pos{\ifgml@bxfirst\global\gml@bxleft\gml@dima\gml@bxfirstfalse\fi
  \global\gml@bxright\gml@dima}
\def\gml@bxrule{\gml@dima\gml@bxright\advance\gml@dima-\gml@bxleft
  \ifdim\gml@dima>\z@\par\nointerlineskip\vskip\gml@rulesep
    \gml@dimb\leftskip\advance\gml@dimb\gml@bxleft
    \moveright\gml@dimb\vbox{\hrule width\gml@dima height\gml@bxwd}\vskip\gml@rulesep\fi}
\gml@bxwd=.3mm
\gmldefhint{bx}{\gml@hget{action}\gml@tv\gml@hget{spec}\gml@tu
  \ifx\gml@tv\gml@offstr
    \ifnum\gml@bxstate=\@ne\gml@bxrule\fi
    \global\gml@bxstate\z@\global\gml@bxleft\z@\global\gml@bxright\z@
  \else\ifx\gml@tv\gml@canstr\global\gml@bxstate\z@
  \else\ifx\gml@tv\gml@charstr
  \else\ifgml@infig\global\gml@bxstate\tw@
  \else
    \ifx\gml@tu\empty\else
      \gml@bxfirsttrue\gml@bxrulefalse
      \global\gml@bxwd=\csname gml@rule@boxrule\endcsname\relax
      \gml@forwords\gml@tu\gml@bx@tok\fi
    \ifdim\gml@bxright=\z@\global\gml@bxright\hsize\fi
    \def\gml@tmpd{line}%
    \ifx\gml@tv\gml@tmpd\ifnum\gml@bxstate=\@ne\gml@bxrule\fi
    \else\ifx\gml@tv\gml@setstr\else\gml@bxrule\fi
      \global\gml@bxstate\@ne\fi
  \fi\fi\fi\fi}
% .hr [rulename] h [FOR w | TO h] ... [LEFT] [RIGHT]: rule segments on
% one line, each in the weight of the last rule name (boxrule).  No
% operands: across the line.  DCF-VERIFY: a start position without FOR,
% TO or RIGHT runs to the right edge of the column.
\def\gml@hr@flush{\gml@dimc\gml@dima\advance\gml@dimc-\gml@hrstart
  \ifdim\gml@dimc>\z@
    \gml@dimd\gml@hrstart\advance\gml@dimd-\gml@hrpos
    \setbox\gml@box\hbox{\unhbox\gml@box\kern\gml@dimd
      \vrule width\gml@dimc height\gml@rulewd depth\z@}%
    \gml@hrpos\gml@dima\fi
  \gml@hrmode\z@}
\def\gml@hr@right{\gml@dima\hsize\advance\gml@dima-\leftskip}
\def\gml@hr@tok#1{\gml@isnumtok{#1}%
  \ifgml@found\gml@dimtok{#1}\gml@dimb
    \ifcase\gml@hrmode \gml@hrstart\gml@dimb\gml@hrmode\@ne
    \or \gml@hr@right\gml@hr@flush\gml@hrstart\gml@dimb\gml@hrmode\@ne
    \or \gml@dima\gml@hrstart\advance\gml@dima\gml@dimb\gml@hr@flush
    \or \gml@dima\gml@dimb\gml@hr@flush\fi
  \else\def\gml@tmpd{#1}%
    \ifx\gml@tmpd\gml@forstr\gml@hrmode\tw@
    \else\ifx\gml@tmpd\gml@tostr\gml@hrmode\thr@@
    \else\ifx\gml@tmpd\gml@leftstr
      \ifnum\gml@hrmode=\@ne\gml@hr@right\gml@hr@flush\fi
      \gml@hrstart-\leftskip\gml@hrmode\@ne
    \else\ifx\gml@tmpd\gml@rightstr
      \ifnum\gml@hrmode=\z@\gml@hrstart-\leftskip\fi
      \gml@hr@right\gml@hr@flush
    \else\ifnum\gml@hrmode=\@ne\gml@hr@right\gml@hr@flush\fi
      \gml@ruleweight{#1}%
    \fi\fi\fi\fi\fi}
\gmldefhint{hr}{\par\gml@hget{spec}\gml@tv
  \gml@rulewd=\csname gml@rule@boxrule\endcsname\relax
  \gml@hrpos\z@\gml@hrmode\z@\setbox\gml@box\hbox{}%
  \ifx\gml@tv\empty\gml@hrstart-\leftskip\gml@hr@right\gml@hr@flush
  \else\gml@forwords\gml@tv\gml@hr@tok
    \ifnum\gml@hrmode=\@ne\gml@hr@right\gml@hr@flush\fi\fi
  \nointerlineskip\vskip\gml@rulesep\moveright\leftskip\box\gml@box
  \vskip\gml@rulesep}
% .kp: keeps cannot be boxed across hints; on/float/delay forbid page
% breaks (between lines and at element skips) until .kp off; .kp v is a
% conditional page eject
\def\gml@kp@tok#1{\gml@dimtok{#1}\gml@dimb\advance\gml@dima\gml@dimb}
\gmldefhint{kp}{\gml@hget{mode}\gml@tv
  \ifx\gml@tv\gml@inlinestr\gml@hget{depth}\gml@tu
    \ifx\gml@tu\empty\else\ifvmode\gml@dima\z@\gml@forwords\gml@tu\gml@kp@tok
      \gml@needspace\gml@dima\fi\fi
  \else\ifx\gml@tv\gml@offstr\global\gml@keepfalse\else\global\gml@keeptrue\fi\fi
  \gml@apply}
\def\gml@runaction#1{\gml@hget{action}\gml@tv
  \ifx\gml@tv\gml@cancelstr\global\csname gml@user#1true\endcsname
    \global\setbox\csname gml@#1odd\endcsname\box\voidb@x
    \global\setbox\csname gml@#1even\endcsname\box\voidb@x
  \else\ifx\gml@tv\gml@supstr\global\csname gml@#1suptrue\endcsname
    \global\csname gml@#1suppagetrue\endcsname
  \else\ifx\gml@tv\gml@resstr\global\csname gml@#1supfalse\endcsname\fi\fi\fi}
\gmldefhint{rh}{\gml@runaction{rh}}
\gmldefhint{rf}{\gml@runaction{rf}}
\gmldefhint{sx}{\gml@hget{left}\gml@sxl\gml@hget{fill}\gml@sxf\gml@hget{right}\gml@sxr
  \par\noindent\gml@sxl
  \ifx\gml@sxf\empty\hfill\else\xleaders\hbox{\gml@sxf}\hfill\fi\gml@sxr\par}
\gmldefhint{tp}{\gml@hget{action}\gml@tv\gml@hget{stops}\gml@tu
  \ifx\gml@tv\gml@resetstr\gml@tabdefault
  \else\ifx\gml@tv\gml@setstr\global\let\gml@tabstops\empty\gml@tabadd
  \else\ifx\gml@tv\gml@addstr\gml@tabadd
  \else\ifx\gml@tv\gml@delstr\gml@tabdel\fi\fi\fi\fi}
\let\gml@hint@tb\gml@hint@tp
% fonts: the first known name of the list (lower case) is used.  An
% unknown list leaves the font unchanged; that is quiet for .bf ... =
% (restart=yes: "the current font is restarted", SH35-0050 p.24) and .pf.
\def\gml@fontlookup#1{\let\gml@fontcmd\relax\gml@lower{#1}%
  \gml@forwords\gml@tmpl\gml@fl@one
  \ifx\gml@fontcmd\relax\ifgml@fontquiet\else
    \gml@warnonce{font#1}{unknown font #1 (font unchanged)}\fi\fi
  \gml@fontquietfalse}
\def\gml@fl@one#1{\ifx\gml@fontcmd\relax
  \expandafter\ifx\csname gml@font@#1\endcsname\relax\else
  \expandafter\let\expandafter\gml@fontcmd\csname gml@font@#1\endcsname\fi\fi}
\gmldefhint{bf}{\gml@hget{font}\gml@tv\gml@hget{restart}\gml@tmpb
  \ifx\gml@tmpb\gml@yes\gml@fontquiettrue\fi\gml@fontlookup\gml@tv
  \ifx\gml@fontcmd\relax\else\gml@hsetfontcmd\fi\gml@apply}
\gmldefhint{pf}{\gml@hget{font}\gml@tv
  \ifx\gml@tv\empty\let\gml@fontcmd\empty\gml@hsetfontcmd
  \else\gml@fontquiettrue\gml@fontlookup\gml@tv
    \ifx\gml@fontcmd\relax\let\gml@fontcmd\empty\fi
    \gml@hsetfontcmd\fi\gml@apply}
% .df NAME spec: tex.xsl parses the spec (os us up font face family size
% bold italic).  A page printer has real faces where a line device
% overstrikes and underscores, so OS (overstrike) is bold and US
% (underscore) is italic: the Starter Set's own page-printer fonts make
% hi1 italic and hi2 bold where its line-device fonts use US and OS
% (DSMPROF3; SH35-0050 p.24-25), and plain TeX cannot underline text by
% a font switch.  UP is done by tex.xsl (\gmlupword).  TYPE: the family
% chooses the slot (mono: tt, sans: ss/ssbf, else rm it bf bi), a size
% loads the slot at that size (\gml@sized).  FONT name alone (no other
% operand) is that font if it is known; DCF-VERIFY: 3800/4250 font ids
% (FONT GT10, &$CHAR(n)) are not interpreted (the body face).  The
% definition is global; \gmldeffont{NAME}{...} after it overrides it.
\gmldefhint{df}{\gml@hget{name}\gml@tv\gml@lower\gml@tv\let\gml@dfname\gml@tmpl
  \gml@hget{os}\gml@dfos\gml@hget{us}\gml@dfus
  \gml@hget{bold}\gml@dfb\gml@hget{italic}\gml@dfi
  \gml@hget{family}\gml@dffam\gml@hget{size}\gml@dfsize\gml@hget{font}\gml@dffont
  \ifx\gml@dfos\gml@yes\let\gml@dfb\gml@yes\fi
  \ifx\gml@dfus\gml@yes\let\gml@dfi\gml@yes\fi
  \def\gml@dfslot{rm}%
  \ifx\gml@dfb\gml@yes\ifx\gml@dfi\gml@yes\def\gml@dfslot{bi}\else\def\gml@dfslot{bf}\fi
  \else\ifx\gml@dfi\gml@yes\def\gml@dfslot{it}\fi\fi
  \ifx\gml@dffam\gml@monostr\def\gml@dfslot{tt}\fi
  \ifx\gml@dffam\gml@sansstr
    \ifx\gml@dfb\gml@yes\def\gml@dfslot{ssbf}\else\def\gml@dfslot{ss}\fi\fi
  \ifx\gml@dfsize\empty\edef\gml@dfcmd{\csname gml@slotcmd@\gml@dfslot\endcsname}%
  \else\edef\gml@dfcmd{\noexpand\gml@sized{\gml@dfslot}{\gml@dfsize}}\fi
  \ifx\gml@dffont\empty\else
    \edef\gml@tmpd{\gml@dfb\gml@dfi\gml@dffam\gml@dfsize}%
    \ifx\gml@tmpd\empty\gml@lower\gml@dffont
      \expandafter\ifx\csname gml@font@\gml@tmpl\endcsname\relax\else
        \expandafter\let\expandafter\gml@dfcmd\csname gml@font@\gml@tmpl\endcsname\fi
    \fi\fi
  \expandafter\global\expandafter\let\csname gml@font@\gml@dfname\endcsname\gml@dfcmd}
\def\gml@slotcmd@rm{\noexpand\rm} \def\gml@slotcmd@it{\noexpand\it}
\def\gml@slotcmd@bf{\noexpand\bf} \def\gml@slotcmd@bi{\noexpand\gml@bi}
\def\gml@slotcmd@tt{\noexpand\tt} \def\gml@slotcmd@ss{\noexpand\gml@ss}
\def\gml@slotcmd@ssbf{\noexpand\gml@ssbf}
% .is [TO] h [MIN m] [BLANK|BREAK|ABSOLUTE] (SH35-0069 p.89-90): h of
% space; TO h, an immediate tab to h, is exact only at the start of a
% line (TeX does not know the position inside a paragraph): elsewhere it
% is the minimum (MIN m, default one word space), as a missed tab.
\def\gml@is@tok#1{\def\gml@tmpd{#1}%
  \ifx\gml@tmpd\gml@tostr\def\gml@isnext{to}%
  \else\ifx\gml@tmpd\gml@minstr\def\gml@isnext{min}%
  \else\gml@isnumtok{#1}\ifgml@found
    \expandafter\let\csname gml@is@\gml@isnext\endcsname\gml@tmpd
    \def\gml@isnext{amount}\fi\fi\fi}
\gmldefhint{is}{\gml@hget{spec}\gml@tv\def\gml@isnext{amount}%
  \let\gml@is@to\empty\let\gml@is@min\empty\let\gml@is@amount\empty
  \gml@forwords\gml@tv\gml@is@tok
  \ifx\gml@is@to\empty
    \ifx\gml@is@amount\empty\else\gml@dimtok\gml@is@amount\gml@dima
      \ifvmode\leavevmode\fi\hskip\gml@dima\fi
  \else\ifvmode\leavevmode\gml@dimtok\gml@is@to\gml@dima\hskip\gml@dima
    \else\ifx\gml@is@min\empty\space
      \else\gml@dimtok\gml@is@min\gml@dima\unskip\hskip\gml@dima\fi\fi\fi}
\gmldefhint{hy}{\gml@hget{mode}\gml@tv
  \ifx\gml@tv\gml@onstr\global\hyphenpenalty50
  \else\ifx\gml@tv\gml@offstr\global\hyphenpenalty\@M\fi\fi}
\def\gml@hw@a#1{\ifx#1\gml@end\let\gml@next\relax\else\let\gml@next\gml@hw@a
  \ifcat a\noexpand#1\edef\gml@hwword{\gml@hwword#1}\else
  \ifx#1-\edef\gml@hwword{\gml@hwword-}\fi\fi\fi\gml@next}
\gmldefhint{hw}{\gml@hget{word}\gml@tv\def\gml@hwword{}%
  \expandafter\gml@hw@a\gml@tv\gml@end
  \ifx\gml@hwword\empty\else\expandafter\hyphenation\expandafter{\gml@hwword}\fi}
\gmldefhint{rc}{\gml@hget{code}\gml@tv\gml@hget{mode}\gml@tu\gml@hget{char}\gml@tw
  \ifx\gml@tv\empty\else
    \ifx\gml@tu\empty\expandafter\global\expandafter\let\csname gml@rc@\gml@tv\endcsname\gml@tw
    \else\ifx\gml@tu\gml@onstr\global\gml@rcontrue
      \expandafter\ifx\csname gml@rc@\gml@tv\endcsname\relax\gdef\gml@rcchar{|}%
      \else\global\expandafter\let\expandafter\gml@rcchar\csname gml@rc@\gml@tv\endcsname\fi
    \else\global\gml@rconfalse\fi\fi\fi}
% .ls normal [v] | by f   (the factor f is kept as written by gml)
\def\gml@twowords#1{\edef\gml@tmpd{#1\space\space}\expandafter\gml@tw@a\gml@tmpd\gml@nil}
\def\gml@tw@a#1 #2 #3\gml@nil{\def\gml@wa{#1}\def\gml@wb{#2}}
\gmldefhint{ls}{\gml@hget{spec}\gml@tv\gml@twowords\gml@tv
  \ifx\gml@wa\gml@bystr
    \ifx\gml@wb\empty\else\global\baselineskip\gml@wb\gml@basebs\fi
  \else\ifx\gml@wa\gml@normalstr
    \ifx\gml@wb\empty\global\baselineskip\gml@basebs
    \else\gml@dimtok\gml@wb\gml@dima\global\baselineskip\gml@dima\fi\fi\fi}
% .ws / .es normal [h] | by f
\def\gml@wordspace#1{\gml@hget{spec}\gml@tv\gml@twowords\gml@tv
  \ifx\gml@wa\gml@bystr\ifx\gml@wb\empty\else
    \global#1=\gml@wb\fontdimen2\font plus\gml@wb\fontdimen3\font
      minus\gml@wb\fontdimen4\font\fi
  \else\ifx\gml@wa\gml@normalstr
    \ifx\gml@wb\empty\global#1\z@skip
    \else\gml@dimtok\gml@wb\gml@dima\global#1\gml@dima plus.5\gml@dima minus.25\gml@dima\fi
  \fi\fi}
\gmldefhint{ws}{\gml@wordspace\spaceskip}
\gmldefhint{es}{\gml@wordspace\xspaceskip}
\def\gml@pn@one#1{\def\gml@tmpd{#1}%
  \ifgml@pnpref\gdef\gml@folioprefix{#1}\gml@pnpreffalse
  \else\ifx\gml@tmpd\gml@prefstr\gml@pnpreftrue
  \else\ifx\gml@tmpd\gml@romanstr\ifnum\pageno>\z@\global\pageno-\pageno\fi
  \else\ifx\gml@tmpd\gml@arabicstr\ifnum\pageno<\z@\global\pageno-\pageno\fi
  \else\ifx\gml@tmpd\gml@offstr\global\gml@nofoliotrue
  \else\ifx\gml@tmpd\gml@offnostr\global\gml@nofoliotrue
  \else\ifx\gml@tmpd\gml@onstr\global\gml@nofoliofalse
  \else\gml@numpart{#1}\ifx\gml@sn\empty\else\global\pageno=\gml@sn\relax\fi
  \fi\fi\fi\fi\fi\fi\fi}
\gmldefhint{pn}{\gml@hget{spec}\gml@tv\gml@pnpreffalse\gml@forwords\gml@tv\gml@pn@one}
% page layout words: the first dimension of the operands
\def\gml@firstdim@one#1{\ifgml@found\else\def\gml@tmpd{#1}%
  \ifx\gml@tmpd\gml@nobindstr\else\gml@dimtok{#1}\gml@dima\gml@foundtrue\fi\fi}
\def\gml@layouthint#1{\gml@hget{spec}\gml@tv\gml@foundfalse
  \gml@forwords\gml@tv\gml@firstdim@one
  \ifgml@found\global#1\gml@dima\gml@layout\fi}
\gmldefhint{pl}{\gml@layouthint\gml@pagelength}
\gmldefhint{pw}{\gml@layouthint\gml@pagewidth}
\gmldefhint{pm}{\gml@layouthint\gml@leftmargin}
\gmldefhint{tm}{\gml@layouthint\gml@topmargin}
\gmldefhint{bm}{\gml@layouthint\gml@botmargin}
% hints with no TeX meaning (or not modelled): a message once
% .cd n positions: the first word is the column count
\gmldefhint{cd}{\gml@hget{spec}\gml@tv\gml@twowords\gml@tv\gml@numpart\gml@wa
  \ifnum0\gml@sn>1 \gml@warnonce{cd}{multi-column layout (.cd) is not supported; one column}\fi}
\gmldefhint{sc}{}
\gmldefhint{mc}{\gml@notimpl{mc}}
\gmldefhint{tc}{\gml@warnonce{tc}{.tc ignored: the table of contents comes from :toc}}
\gmldefhint{vr}{\gml@notimpl{vr}}
\gmldefhint{sb}{\gml@notimpl{sb}}
\gmldefhint{bc}{} \gmldefhint{fl}{} \gmldefhint{wz}{}
% .sa/.re: the engine follows a .re with the hints that set the saved
% state again (restore=yes), so nothing to do here.  .fv: page-level
% vertical justification is TeX's own business.
\gmldefhint{sa}{} \gmldefhint{re}{} \gmldefhint{fv}{}
\gmldefhint{ic}{\gml@notimpl{ic}}
\gmldefhint{ud}{\gml@notimpl{ud}}

% ---- line scopes ------------------------------------------------------------------------
% .ce/.ri: each covered input line is a centred / right-aligned line
\def\gml@linesalign#1{\def\gml@curfill{no}\def\gml@curalign{#1}%
  \gml@applymargins\catcode`\^^M=\active}
\gmldeflines{ce}{begin}{\gml@linesalign{center}}
\gmldeflines{ce}{end}{\gml@lmore}
\gmldeflines{ri}{begin}{\gml@linesalign{right}}
\gmldeflines{ri}{end}{\gml@lmore}
% .us/.up/.uc: tex.xsl wraps the words (\gmlusword ...); nothing else
\gmldeflines{us}{begin}{} \gmldeflines{up}{begin}{} \gmldeflines{uc}{begin}{}
% .in/.ir FOR/AFTER n, .il, .un: the first n lines get the scope's margin.
% Fill mode: \parshape over n output lines (set before the group);
% format off: the margin for the covered lines (each one a paragraph).
\def\gml@lines@margin#1#2#3{% attr, permanent register, left|right
  \gml@get\gml@lattrs{#1}\gml@tv\gml@get\gml@lattrs{unit}\gml@tu
  \gml@dimof\gml@tv\gml@tu\gml@dima\advance\gml@dima-#2%
  \edef\gml@tdelta{\the\gml@dima}%
  \gml@get\gml@lattrs{n}\gml@tv\gml@numpart\gml@tv
  \ifx\gml@sn\empty\def\gml@sn{1}\fi
  \ifx\gml@curfill\gml@no\else
    \gml@get\gml@lattrs{cont}\gml@tmpb\ifx\gml@tmpb\gml@yes\else
      \def\gml@tmpb{left}\def\gml@tmpc{#3}%
      \ifx\gml@tmpb\gml@tmpc\gml@shape0\gml@sn\gml@tdelta\z@
      \else\gml@shape0\gml@sn\z@\gml@tdelta\fi\fi\fi}
\gmldeflines{in}{pre}{\gml@lines@margin{left}\gml@hintleft{left}}
\gmldeflines{in}{begin}{\ifx\gml@curfill\gml@no
  \advance\leftskip\gml@tdelta\relax\fi}
\gmldeflines{in}{end}{\ifx\gml@curfill\gml@no\gml@lmore\fi}
\gmldeflines{ir}{pre}{\gml@lines@margin{right}\gml@hintright{right}}
\gmldeflines{ir}{begin}{\ifx\gml@curfill\gml@no\advance\rightskip\gml@tdelta\relax\fi}
\gmldeflines{ir}{end}{\ifx\gml@curfill\gml@no\gml@lmore\fi}
\expandafter\let\csname gml@lines@il@pre\expandafter\endcsname\csname gml@lines@in@pre\endcsname
\expandafter\let\csname gml@lines@il@begin\expandafter\endcsname\csname gml@lines@in@begin\endcsname
\expandafter\let\csname gml@lines@il@end\expandafter\endcsname\csname gml@lines@in@end\endcsname
\expandafter\let\csname gml@lines@un@pre\expandafter\endcsname\csname gml@lines@in@pre\endcsname
\expandafter\let\csname gml@lines@un@begin\expandafter\endcsname\csname gml@lines@in@begin\endcsname
\expandafter\let\csname gml@lines@un@end\expandafter\endcsname\csname gml@lines@in@end\endcsname
% .rc * c: a revision character in the margin of the covered line
\gmldeflines{rc}{begin}{\gml@lget{char}\gml@tv\gml@lget{cont}\gml@tmpb
  \ifx\gml@tmpb\gml@yes\else\leavevmode\llap{\gml@tv\enspace}\fi}

% ---- page references ----------------------------------------------------------------
\def\gml@pagenumber#1{\expandafter\ifx\csname gml@pg@#1\endcsname\relax
  \global\gml@refschangedtrue??\else\csname gml@pg@#1\endcsname\fi}
\def\gmlanchor#1{\ifgml@reffile
  \edef\gml@tmpa{\write\gml@refout{\string\gml@ap{#1}{\noexpand\gml@folio}}}%
  \ifvmode\gml@tmpa\nobreak\else\gml@tmpa\fi\fi}
\def\gml@reflabel#1{\def\gml@tmpe{#1}\ifx\gml@tmpe\empty\else\ #1\nobreak\ \fi}
\def\gmlpageref#1#2#3{\def\gml@tmpd{#2}%
  \ifx\gml@tmpd\gml@autostr\gml@pageref@auto{#1}{#3}%
  \else\gml@reflabel{#3}\gml@pagenumber{#1}\fi}
\def\gml@pageref@auto#1#2{\global\advance\gml@refcount\@ne
  \ifgml@reffile\edef\gml@tmpa{\write\gml@refout
    {\string\gml@rp{\the\gml@refcount}{\noexpand\gml@folio}}}\gml@tmpa\fi
  \expandafter\let\expandafter\gml@pa\csname gml@rpg@\the\gml@refcount\endcsname
  \expandafter\let\expandafter\gml@pb\csname gml@pg@#1\endcsname
  \ifx\gml@pa\relax\gml@reflabel{#2}\gml@pagenumber{#1}%
  \else\ifx\gml@pa\gml@pb\else\gml@reflabel{#2}\gml@pagenumber{#1}\fi\fi}
\def\gml@folio{\gml@folioprefix\folio}
\def\gml@readrefs{\begingroup\catcode`\@=11
  \def\gml@ap##1##2{\expandafter\gdef\csname gml@pg@##1\endcsname{##2}}%
  \def\gml@rp##1##2{\expandafter\gdef\csname gml@rpg@##1\endcsname{##2}}%
  \openin\gml@in=\jobname.ref
  \ifeof\gml@in\closein\gml@in\global\gml@refschangedtrue
  \else\closein\gml@in\input\jobname.ref \fi\endgroup
  \immediate\openout\gml@refout=\jobname.ref \global\gml@reffiletrue}
\def\gml@checkrefs{\immediate\closeout\gml@refout\global\gml@reffilefalse
  \begingroup\catcode`\@=11
  \def\gml@ap##1##2{\def\gml@tmpd{##2}\expandafter\ifx\csname gml@pg@##1\endcsname
    \gml@tmpd\else\global\gml@refschangedtrue\fi}%
  \def\gml@rp##1##2{\def\gml@tmpd{##2}\expandafter\ifx\csname gml@rpg@##1\endcsname
    \gml@tmpd\else\global\gml@refschangedtrue\fi}%
  \openin\gml@in=\jobname.ref
  \ifeof\gml@in\closein\gml@in\else\closein\gml@in\input\jobname.ref \fi\endgroup
  \ifgml@refschanged\immediate\write16{GML: page references have changed;
    run TeX again to get them right.}\fi}

% ---- running headings and footings ------------------------------------------------------
\def\gml@sec{} \def\gml@shead{}
\headline={\gml@headline} \footline={\gml@footline}
% .rh/.rf sup suppresses the page it is issued on even if res follows
% before the page is complete (the Starter Set's title page: sup, the
% title page text, res, then the page eject)
\def\gml@headline{\rm\ifgml@plainpage\hfil\else\ifgml@rhsuppage\hfil
  \else\ifgml@userrh\gml@userrh\else\gml@defaultrh\fi\fi\fi}
\def\gml@footline{\rm\ifgml@plainpage\hfil\else\ifgml@rfsuppage\hfil
  \else\ifgml@userrf\gml@userrf\else\gml@defaultrf\fi\fi\fi}
% Starter Set: the security classification centred at the top; the short
% heading on the left and the page number on the right at the foot
\def\gml@defaultrh{\hfil{\gml@hi{2}\gml@sec}\hfil}
\def\gml@defaultrf{\gml@shead\hfil\ifgml@nofolio\else\gml@folio\fi}
\def\gml@userrh{\ifodd\pageno\copy\gml@rhodd\else\copy\gml@rheven\fi\hss}
\def\gml@userrf{\ifodd\pageno\copy\gml@rfodd\else\copy\gml@rfeven\fi\hss}
\output={\hsize\gml@defaulthsize\plainoutput\global\gml@plainpagefalse
  \global\gml@chappagefalse
  \ifgml@rhsup\global\gml@rhsuppagetrue\else\global\gml@rhsuppagefalse\fi
  \ifgml@rfsup\global\gml@rfsuppagetrue\else\global\gml@rfsuppagefalse\fi}
\def\gml@runsetup{\hsize\gml@defaulthsize\gml@elemhsize\hsize
  \gml@elemleft\z@\gml@elemright\z@\def\gml@efont{\rm}\gml@apply}
% a gml:rh/gml:rf with action= (.rh sup|res|cancel as an element) only
% acts; it has no text
\def\gml@runstart#1#2{\gml@aget{page}\gml@runpage\gml@aget{action}\gml@runact
  \ifx\gml@runact\empty\else\let\gml@hattrs\gml@attrs\gml@runaction{#1}\fi
  \setbox\gml@box#2\bgroup\gml@runsetup}
\def\gml@runend#1{\par\egroup\ifx\gml@runact\empty\gml@runstore{#1}\fi}
\def\gml@runstore#1{\global\csname gml@user#1true\endcsname
  \ifx\gml@runpage\gml@evenstr\else
    \global\setbox\csname gml@#1odd\endcsname\copy\gml@box\fi
  \ifx\gml@runpage\gml@oddstr\else
    \global\setbox\csname gml@#1even\endcsname\copy\gml@box\fi}
\gmldefstart{gml:rh}{\gml@runstart{rh}\vbox}
\gmldefend{gml:rh}{\gml@runend{rh}}
\gmldefstart{gml:rf}{\gml@runstart{rf}\vtop}
\gmldefend{gml:rf}{\gml@runend{rf}}

% ---- headings ------------------------------------------------------------------------------
% SH35-0050 Figure 6 (fonts) and Figure 14 (spacing, page devices).
% new@N: a new page before; sp@N: the space before; post@N: after.
\expandafter\let\csname gml@head@new@0\endcsname\gml@yes
\expandafter\let\csname gml@head@new@1\endcsname\gml@yes
\expandafter\def\csname gml@head@sp@0\endcsname{\gml@vspace{1.3in}}
\expandafter\def\csname gml@head@sp@1\endcsname{\gml@vspace{1.3in}}
\expandafter\def\csname gml@head@sp@2\endcsname{\gml@vskip{20bp}}
\expandafter\def\csname gml@head@sp@3\endcsname{\gml@vskip{18bp}}
\expandafter\def\csname gml@head@sp@4\endcsname{\gml@vskip{14bp}}
\expandafter\def\csname gml@head@sp@5\endcsname{\gml@vskip\baselineskip}
\expandafter\def\csname gml@head@sp@6\endcsname{\gml@vskip\baselineskip}
\expandafter\def\csname gml@head@post@0\endcsname{\gml@vskip{14bp}}
\expandafter\def\csname gml@head@post@1\endcsname{\gml@vskip{14bp}}
\expandafter\def\csname gml@head@post@2\endcsname{\gml@vskip{11bp}}
\expandafter\def\csname gml@head@post@3\endcsname{\gml@vskip{11bp}}
\expandafter\def\csname gml@head@post@4\endcsname{\gml@vskip{11bp}}
\expandafter\def\csname gml@head@post@5\endcsname{}
\expandafter\def\csname gml@head@post@6\endcsname{}
\expandafter\def\csname gml@head@bs@0\endcsname{25pt}
\expandafter\def\csname gml@head@bs@1\endcsname{25pt}
\expandafter\def\csname gml@head@bs@2\endcsname{21pt}
\expandafter\def\csname gml@head@bs@3\endcsname{17pt}
\expandafter\def\csname gml@head@bs@4\endcsname{14.5pt}
\expandafter\def\csname gml@head@bs@5\endcsname{12pt}
\expandafter\def\csname gml@head@bs@6\endcsname{12pt}
% .dh n options (tex.xsl parses them): \gml@dh@N@KEY overrides the
% defaults above, for headings of level N (:hN and .hN alike) until a
% bare .dh n.  font (a .df/known font; its baseline is at least the body
% one), pa=yes|no, spbf (a space, kept at a page top), skbf (a skip),
% spaf (after; a skip), align, runin=yes (NBR: the text runs on), hy.
% UP is done by tex.xsl; US, NUM, HANG, SECT and the table of contents
% options are not used (DCF-VERIFY: US might be meant as italic).
\gmldefhint{dh}{\gml@hget{level}\gml@dhl\gml@hget{reset}\gml@tv
  \ifx\gml@tv\gml@yes\gml@forwords{font pa spbf skbf spaf align runin hy}\gml@dh@clear\fi
  \gml@forwords{font pa spbf skbf spaf align runin hy}\gml@dh@set}
\def\gml@dh@set#1{\gml@hget{#1}\gml@tv\ifx\gml@tv\empty\else
  \expandafter\global\expandafter\let\csname gml@dh@\gml@dhl @#1\endcsname\gml@tv\fi}
\def\gml@dh@clear#1{\expandafter\global\expandafter\let
  \csname gml@dh@\gml@dhl @#1\endcsname\relax}
\def\gml@dhget#1#2{\expandafter\let\expandafter\gml@dhv\csname gml@dh@#1@#2\endcsname
  \ifx\gml@dhv\relax\let\gml@dhv\empty\fi}
\def\gml@headstart#1{\par
  \gml@dhget{#1}{pa}\let\gml@dhpa\gml@dhv
  \ifx\gml@dhpa\empty\expandafter\let\expandafter\gml@dhpa\csname gml@head@new@#1\endcsname\fi
  \ifx\gml@dhpa\gml@yes\gml@chapnewpage\fi
  \gml@dhget{#1}{spbf}\let\gml@dhspbf\gml@dhv\gml@dhget{#1}{skbf}\let\gml@dhskbf\gml@dhv
  \ifx\gml@dhspbf\empty\ifx\gml@dhskbf\empty\csname gml@head@sp@#1\endcsname\fi
  \else\gml@dimtok\gml@dhspbf\gml@dima\gml@vspace\gml@dima\fi
  \ifx\gml@dhskbf\empty\else\gml@dimtok\gml@dhskbf\gml@dima\gml@vskip\gml@dima\fi
  \gml@dhget{#1}{font}\let\gml@fontcmd\relax
  \ifx\gml@dhv\empty\else\gml@fontlookup\gml@dhv\fi
  \ifx\gml@fontcmd\relax
    \gml@setfont{\csname gml@font@hd#1\endcsname}%
    \baselineskip\csname gml@head@bs@#1\endcsname\relax
  \else\expandafter\gml@setfont\expandafter{\gml@fontcmd}%
    \gml@dima=1.2\fontdimen6\font
    \ifdim\gml@dima<\gml@basebs\gml@dima\gml@basebs\fi\baselineskip\gml@dima\fi
  \gml@dhget{#1}{hy}\ifx\gml@dhv\gml@yes\hyphenpenalty50 \else\hyphenpenalty\@M\fi
  \gml@dhget{#1}{align}\ifx\gml@dhv\empty\def\gml@dhv{left}\fi
  \gml@setfill{yes}\gml@dhv\gml@apply\gml@inheadtrue\noindent}
\def\gml@headend#1{\gml@dhget{#1}{runin}%
  \ifx\gml@dhv\gml@yes\unskip\enspace\kern1sp
  \else\par\nobreak
    \gml@dhget{#1}{spaf}\ifx\gml@dhv\empty\csname gml@head@post@#1\endcsname
    \else\gml@dimtok\gml@dhv\gml@dima\gml@vskip\gml@dima\fi\fi}
\def\gml@headshead{\gml@aget{stitle}\gml@tv\ifx\gml@tv\empty\gml@aget{text}\gml@tv\fi
  \global\let\gml@shead\gml@tv}
\def\gml@headnum#1{#1\quad}
\def\gml@defheading#1{%
  \gmldefstart{h#1}{\gml@headstart{#1}\ifnum#1<2 \gml@headshead\fi}%
  \gmldefend{h#1}{\gml@headend{#1}}%
  \expandafter\let\csname gml@num@h#1\endcsname\gml@headnum}
\gml@defheading0 \gml@defheading1 \gml@defheading2 \gml@defheading3
\gml@defheading4 \gml@defheading5 \gml@defheading6
\gmldefstart{gml:h}{\gml@aget{level}\gml@lvl\ifx\gml@lvl\empty\def\gml@lvl{1}\fi
  \expandafter\gml@headstart\expandafter{\gml@lvl}}
\gmldefend{gml:h}{\expandafter\gml@headend\expandafter{\gml@lvl}}
% a heading with a literal text (abstract, preface, toc, figlist, index)
\def\gml@sectionhead#1{\begingroup\gml@headstart1#1\gml@headend1\endgroup
  \gdef\gml@shead{#1}\gml@apply}

% ---- document sections and the title page -----------------------------------------------
\gmldefstart{gdoc}{\gml@aget{sec}\gml@tv\global\let\gml@sec\gml@tv}
\gmldefend{gdoc}{\par}
\gmldefstart{frontm}{\gml@newpage\global\pageno\m@ne}
\gmldefstart{body}{\gml@newpage\global\pageno\@ne}
\gmldefstart{appendix}{\gml@newpage}
\gmldefstart{backm}{\gml@newpage}
\gmldefstart{abstract}{\gml@newpage\gml@sectionhead\gmlabstracttitle}
\gmldefstart{preface}{\gml@newpage\gml@sectionhead\gmlprefacetitle}
\gmldefstart{toc}{} \gmldefstart{figlist}{} \gmldefstart{index}{}
\gmldefstart{titlep}{\gml@newpage\global\gml@plainpagetrue\gml@intitleptrue
  \gml@vspace\gml@titletop\global\gml@aftertitlefalse
  \gml@setfill{yes}{\gmltitlealign}\hyphenpenalty\@M\gml@apply}
\gmldefend{titlep}{\par\ifx\gml@sec\empty\else\gml@vskip{2\baselineskip}%
  \noindent{\gml@hi{2}\gml@sec}\par\fi\gml@newpage}
\def\gml@titlegap{\par\ifgml@aftertitle\global\gml@aftertitlefalse
  \gml@vspace\gml@titlegapdim\else\gml@vskip\baselineskip\fi}
\gmldefstart{title}{\par\gml@setfont{\csname gml@font@title\endcsname}%
  \baselineskip30pt\gml@apply
  \gml@aget{stitle}\gml@tv\ifx\gml@tv\empty\gml@aget{text}\gml@tv\fi
  \global\let\gml@shead\gml@tv\noindent}
\gmldefend{title}{\par\global\gml@aftertitletrue}
\gmldefstart{docnum}{\gml@titlegap\noindent\gmldocnumlabel\ }
\gmldefend{docnum}{\par}
\gmldefstart{date}{\gml@titlegap\noindent}
\gmldefend{date}{\par}
\gmldefstart{author}{\gml@titlegap\noindent}
\gmldefend{author}{\par}
\gmldefstart{address}{\ifgml@intitlep\gml@titlegap\gml@setfill{no}{\gmltitlealign}%
  \else\gml@vskip{.75\baselineskip}\advance\gml@elemleft4\gml@chwidth
    \gml@setfill{no}{left}\fi\gml@apply}
\gmldefend{address}{\par\ifgml@intitlep\else\gml@vskip{.75\baselineskip}\fi}
\gmldefstart{aline}{\par}
\gmldefend{aline}{\par}

% ---- paragraphs, notes, quotations, examples --------------------------------------------
\gmldefstart{p}{\gml@parstart{.75\baselineskip}}
\gmldefend{p}{\par}
\gml@letelement{pc}{p}
\gmldefstart{note}{\gml@parstart{.75\baselineskip}\leavevmode
  {\gml@hi{2}\gmlnotelabel}\nobreak\ }
\gmldefend{note}{\par}
\gmldefstart{lq}{\gml@vskip\baselineskip\advance\gml@elemleft3\gml@chwidth
  \advance\gml@elemright3\gml@chwidth\gml@apply}
\gmldefend{lq}{\gml@vskip\baselineskip}
\gmldefstart{xmp}{\gml@vskip\baselineskip\advance\gml@elemleft2\gml@chwidth
  \gml@setfont\tt\gml@setfill{no}{left}\gml@preservetrue\gml@apply}
\gmldefend{xmp}{\par\gml@vskip\baselineskip}

% ---- highlighting and phrases -------------------------------------------------------------
% in a heading the heading font is kept (it is already bold and large)
\def\gml@hp#1{\ifgml@inhead\else\gml@setfont{\gml@hi{#1}}\fi}
\gmldefstart{hp0}{\gml@hp0}
\gmldefstart{hp1}{\gml@hp1} \gmldefend{hp1}{\/}
\gmldefstart{hp2}{\gml@hp2}
\gmldefstart{hp3}{\gml@hp3} \gmldefend{hp3}{\/}
\gmldefstart{cit}{\gml@setfont{\gml@hi{1}}} \gmldefend{cit}{\/}
\gmldefstart{q}{\advance\gml@qdepth\@ne\ifodd\gml@qdepth``\else`\fi}
\gmldefend{q}{\ifodd\gml@qdepth''\else'\fi}
\gmldefstart{hdref}{} \gmldefstart{figref}{} \gmldefstart{liref}{}
\def\gml@fnmark#1{\leavevmode\raise.5ex\hbox{\sevenrm#1}}
\gmldefstart{fnref}{\leavevmode\raise.5ex\hbox\bgroup\sevenrm}
\gmldefend{fnref}{\egroup}
% index tags: only their anchors (the terms are in gml:index-entries)
\gmldefstart{i1}{} \gmldefstart{i2}{} \gmldefstart{i3}{}
\gmldefstart{ih1}{} \gmldefstart{ih2}{} \gmldefstart{ih3}{} \gmldefstart{iref}{}

% ---- lists ---------------------------------------------------------------------------------------
% SH35-0050 Figure 5: indent / skip  ol ul sl 4 .75, dl 10 1, gl 0 .75
\def\gml@itemskip{\z@} \def\gml@listtype{} \def\gml@ulnest{1}
\def\gml@termhi{2} \def\gml@headhi{3}
% \gml@itembreak: no page break between the first and second item of a list
% (a club line: a lone first item at the foot of a page)
\def\gml@itembreak{\global\advance\gml@itemno\@ne
  \ifnum\gml@itemno=\tw@ \par\nobreak\fi}
\def\gml@outeritem{0}
\def\gml@liststart#1#2#3{\gml@vskip{#3}\def\gml@listtype{#1}%
  \edef\gml@outeritem{\the\gml@itemno}\global\gml@itemno\z@
  \advance\gml@elemleft#2\relax\gml@listind#2\relax
  \gml@aget{compact}\gml@tv\ifx\gml@tv\empty\def\gml@itemskip{#3}\else\def\gml@itemskip{\z@}\fi
  \gml@aget{termhi}\gml@tv\ifx\gml@tv\empty\else\let\gml@termhi\gml@tv\fi
  \gml@aget{headhi}\gml@tv\ifx\gml@tv\empty\else\let\gml@headhi\gml@tv\fi
  \gml@apply}
\def\gml@listend#1{\global\gml@itemno\gml@outeritem\relax\gml@vskip{#1}}
\gmldefstart{ol}{\gml@liststart{ol}{4\gml@chwidth}{.75\baselineskip}}
\gmldefend{ol}{\gml@listend{.75\baselineskip}}
\gml@letelement{ul}{ol}
\gmldefstart{ul}{\gml@liststart{ul}{4\gml@chwidth}{.75\baselineskip}%
  \gml@aget{nest}\gml@tv\ifx\gml@tv\empty\else\let\gml@ulnest\gml@tv\fi}
\gmldefstart{sl}{\gml@liststart{sl}{4\gml@chwidth}{.75\baselineskip}}
\gmldefend{sl}{\gml@listend{.75\baselineskip}}
\gml@letelement{l}{sl}
\gmldefstart{dl}{\gml@liststart{dl}{10\gml@chwidth}{\baselineskip}%
  \gml@aget{break}\gml@dlbreak}
\gmldefend{dl}{\gml@listend\baselineskip}
\gmldefstart{gl}{\gml@liststart{gl}{\z@}{.75\baselineskip}}
\gmldefend{gl}{\gml@listend{.75\baselineskip}}
\def\gmlbullet#1{\ifcase#1\or$\bullet$\or---\or$\ast$\else$\bullet$\fi}
\def\gml@itemlabel#1{\llap{\hbox to\gml@listind{#1\hss}}}
\def\gml@li@ul{\gml@itemlabel{\gmlbullet\gml@ulnest}}
\gmldefstart{li}{\par\gml@itembreak\gml@vskip\gml@itemskip\leavevmode
  \csname gml@li@\gml@listtype\endcsname}
\gmldefend{li}{\par}
\gmldefnum{li}#1{\gml@itemlabel{#1}}
\gmldefstart{lp}{\par\gml@vskip\gml@itemskip}
\gmldefend{lp}{\par}
% terms: a box in the hanging indent, or a line of its own if too wide
\def\gml@termstart#1{\par\gml@itembreak\gml@vskip\gml@itemskip\noindent\hskip-\gml@listind
  \setbox\gml@dtbox\hbox\bgroup\gml@setfont{\gml@hi{#1}}}
\def\gml@termend{\unskip\egroup\gml@dimb\gml@listind\advance\gml@dimb-\gml@chwidth
  \ifx\gml@dlbreak\empty\else\gml@dimb\z@\fi
  \ifdim\wd\gml@dtbox>\gml@dimb\box\gml@dtbox\par
  \else\hbox to\gml@listind{\box\gml@dtbox\hss}\kern1sp\fi}
\gmldefstart{dt}{\gml@termstart\gml@termhi}
\gmldefend{dt}{\gml@termend}
\gmldefstart{dthd}{\gml@termstart\gml@headhi}
\gmldefend{dthd}{\gml@termend}
\gmldefstart{dd}{}
\gmldefend{dd}{\par}
\gmldefstart{ddhd}{\gml@setfont{\gml@hi{\gml@headhi}}}
\gmldefend{ddhd}{\par}
\let\gml@dlbreak\empty
\gmldefstart{gt}{\par\gml@itembreak\gml@vskip\gml@itemskip\noindent\gml@setfont{\gml@hi{\gml@termhi}}}
\gmldefend{gt}{\unskip\/\kern1sp}
\gmldefstart{gd}{\ifhmode\ifdim\lastkern=1sp\unkern:\enspace\fi\fi}
\gmldefend{gd}{\par}

% ---- figures ---------------------------------------------------------------------------------------
% frame=rule (default): rules above and below; box; none; any other value:
% a line of that string.  place=inline: \midinsert, else \topinsert (plain
% TeX has no bottom floats).  On a page that begins with an h0/h1 or
% section heading a top float would go above that heading, so there it is
% a \midinsert too: here if it fits, else at the top of the next page.
% The body is in format-off mode.
\def\gml@charrule#1{\hbox to\hsize{\xleaders\hbox{#1}\hfil}}
\def\gml@figopen{\par\gml@aget{place}\gml@tv
  \ifx\gml@tv\gml@inlinestr\midinsert
  \else\ifgml@chappage\midinsert\else\topinsert\fi\fi
  \gml@aget{frame}\gml@figframe\ifx\gml@figframe\empty\let\gml@figframe\gml@rulestr\fi
  \setbox\gml@figbox\vbox\bgroup\gml@infigtrue
  \ifx\gml@figframe\gml@boxstr\advance\gml@elemhsize-2\gml@figpad
    \advance\gml@elemhsize-.8pt\fi}
\gmldefstart{fig}{\gml@figopen
  \advance\gml@elemleft2\gml@chwidth\advance\gml@elemright2\gml@chwidth
  \gml@setfill{no}{left}\gml@apply}
\gmldefhstart{fig}{\gml@figopen\gml@apply}
\gmldefhend{fig}{\csname gml@fig@end\endcsname}
\gmldefend{fig}{\par\egroup
  \ifx\gml@figframe\gml@rulestr\hrule\kern\gml@figpad\box\gml@figbox\kern\gml@figpad\hrule
  \else\ifx\gml@figframe\gml@boxstr\vbox{\hrule\hbox{\vrule\kern\gml@figpad
      \vbox{\kern\gml@figpad\box\gml@figbox\kern\gml@figpad}\kern\gml@figpad\vrule}\hrule}%
  \else\ifx\gml@figframe\gml@nonestr\box\gml@figbox
  \else\gml@charrule\gml@figframe\box\gml@figbox\gml@charrule\gml@figframe
  \fi\fi\fi\endinsert}
\gmldefstart{figcap}{\gml@vskip\baselineskip\gml@setfill{yes}{left}\gml@apply\noindent}
\gmldefend{figcap}{\gml@aget{desc}\gml@tv\ifx\gml@tv\gml@yes\unskip\kern1sp\else\par\fi}
\gmldefnum{figcap}#1{\setbox\gml@box\hbox{{\bf#1}\enspace}\global\gml@capind\wd\gml@box
  \hangindent\gml@capind\hangafter\@ne\box\gml@box}
\gmldefstart{figdesc}{\gml@setfill{yes}{left}\gml@apply
  \ifhmode\ifdim\lastkern=1sp\unkern\fi:\ \ \hangindent\gml@capind\hangafter\@ne
  \else\gml@vskip\baselineskip\noindent\fi}
\gmldefend{figdesc}{\par}

% ---- footnotes -------------------------------------------------------------------------------------------
\def\gml@vfootnote#1{\insert\footins\bgroup
  \interlinepenalty\interfootnotelinepenalty
  \splittopskip\ht\strutbox\splitmaxdepth\dp\strutbox\floatingpenalty\@MM
  \leftskip\z@skip\rightskip\z@skip\spaceskip\z@skip\xspaceskip\z@skip
  \parindent1.5em\textindent{#1}\footstrut\futurelet\next\fo@t}
\def\gml@footnote#1{\let\@sf\empty\ifhmode\edef\@sf{\spacefactor\the\spacefactor}\/\fi
  #1\@sf\gml@vfootnote{#1}}
\def\gml@fnsetup{\gml@elemleft\z@\gml@elemright\z@\gml@elemhsize\hsize}
\gmldefstart{fn}{\gml@aget{num}\gml@tv\gml@aget{callout}\gml@tmpb
  \edef\gml@tmpc{\noexpand\gml@fnmark{\gml@tv}}%
  \ifx\gml@tmpb\gml@no\let\gml@fncmd\gml@vfootnote\else\let\gml@fncmd\gml@footnote\fi
  \expandafter\gml@fncmd\expandafter{\gml@tmpc}\bgroup\gml@fnsetup}
\gmldefend{fn}{\egroup}
\gmldefnum{fn}#1{}
\gmldefstart{gml:fn}{\gml@aget{kind}\gml@fnkind
  \ifx\gml@fnkind\gml@leaderstr\global\gml@userleadertrue
    \global\setbox\gml@fnleader\vbox\bgroup\gml@runsetup
  \else\gml@vfootnote{}\bgroup\gml@fnsetup\fi}
\gmldefend{gml:fn}{\ifx\gml@fnkind\gml@leaderstr\par\egroup\else\egroup\fi}
\let\gml@plainfootnoterule\footnoterule
\def\footnoterule{\ifgml@userleader\copy\gml@fnleader\kern2pt
  \else\gml@plainfootnoterule\fi}

% ---- table of contents, list of illustrations ---------------------------------------------------
\def\gml@tocind#1{\ifnum#1<2 \gml@dima\z@\else\gml@cnta#1\advance\gml@cnta\m@ne
  \multiply\gml@cnta\tw@\gml@dima\gml@cnta\gml@chwidth\fi}
\def\gml@leaderline#1#2#3{% num text page
  \leftskip\gml@dima\rightskip\gml@tocpage plus1fil\parfillskip-\rightskip
  \def\gml@tmpd{#1}\ifx\gml@tmpd\empty\else
    \setbox\gml@box\hbox{#1\enspace}\hangindent\wd\gml@box\hangafter\@ne\fi
  \noindent\ifx\gml@tmpd\empty\else\box\gml@box\fi #2\nobreak
  \leaders\hbox to.75em{\hss.\hss}\hfill\nobreak\hbox to\gml@tocpage{\hss#3}\par}
\def\gmltocbegin#1{\gml@sectionhead{#1}\begingroup\hyphenpenalty\@M}
\def\gmltocentry#1#2#3#4{\def\gml@tattrs{#1}\gml@get\gml@tattrs{level}\gml@tv
  \ifx\gml@tv\empty\def\gml@tv{1}\fi
  \par\ifnum\gml@tv<2 \gml@vskip\baselineskip\fi
  \begingroup\gml@tocind\gml@tv\advance\gml@dima\gml@elemleft
  \ifnum\gml@tv<2 \bf\fi\gml@leaderline{#2}{#3}{#4}\endgroup}
\def\gmltocend{\par\endgroup\gml@apply\gml@newpage}
\let\gmlfiglistbegin\gmltocbegin
\def\gmlfiglistentry#1#2#3#4{\par\gml@vskip{.5\baselineskip}%
  \begingroup\gml@dima\gml@elemleft\gml@leaderline{#2}{#3}{#4}\endgroup}
\let\gmlfiglistend\gmltocend

% ---- index -----------------------------------------------------------------------------------------------
% \gmlindexcolumns = 1 or 2 (default 2, as DSMINDEX); two columns are
% balanced with \vsplit.
\def\gmlindexbegin#1{\gml@sectionhead{#1}\begingroup\hyphenpenalty\@M
  \ifnum\gmlindexcolumns>\@ne
    \gml@dima\hsize\advance\gml@dima-\gml@colsep\divide\gml@dima\tw@
    \setbox\gml@ixbox\vbox\bgroup\hsize\gml@dima\gml@elemhsize\gml@dima\gml@apply
  \fi}
\def\gmlindexend{\par\ifnum\gmlindexcolumns>\@ne\egroup\gml@ixsplit\fi\endgroup\gml@apply}
\def\gml@ixsplit{\begingroup\splittopskip\baselineskip\splitmaxdepth\maxdepth
  \vbadness\@M\gml@cntb\z@
  \loop\ifdim\ht\gml@ixbox>\z@\advance\gml@cntb\@ne
    \endgraf\penalty\@M
    \ifdim\pagegoal=\maxdimen\gml@dima\vsize
    \else\gml@dima\pagegoal\advance\gml@dima-\pagetotal\advance\gml@dima-\baselineskip\fi
    \ifdim\gml@dima<2\baselineskip\vfill\penalty-\@M\gml@dima\vsize\fi
    \gml@dimb\ht\gml@ixbox\advance\gml@dimb\dp\gml@ixbox\divide\gml@dimb\tw@
    \advance\gml@dimb\baselineskip
    \ifdim\gml@dimb>\gml@dima\gml@dimb\gml@dima\fi
    \setbox\z@\vsplit\gml@ixbox to\gml@dimb
    \setbox\tw@\vsplit\gml@ixbox to\gml@dimb
    \hbox to\hsize{\vbox to\gml@dimb{\unvbox\z@\vss}\hfil\vbox to\gml@dimb{\unvbox\tw@\vss}}%
    \ifnum\gml@cntb>500 \setbox\gml@ixbox\box\voidb@x\fi
  \repeat\endgroup}
\def\gmlindexgroup#1{\par\gml@vskip\baselineskip\def\gml@tmpd{#1}%
  \ifx\gml@tmpd\empty\else\noindent{\bf#1}\par\nobreak\fi}
\def\gmlindexentry#1#2#3{\par\begingroup\gml@cnta#1\relax\advance\gml@cnta\m@ne
  \gml@dima\gml@cnta\gml@chwidth\multiply\gml@dima\tw@\advance\leftskip\gml@dima
  \rightskip\z@ plus3em\hangindent4\gml@chwidth\hangafter\@ne
  \noindent#2#3\par\endgroup}
\def\gmlindexfirst{,\ }
\def\gmlindexsep{,\ }
\def\gmlindexrangesep{--}
\def\gmlindexref#1#2{\def\gml@xattrs{#2}\gml@get\gml@xattrs{text}\gml@tv
  \gml@get\gml@xattrs{major}\gml@tu
  \ifx\gml@tv\empty\ifx\gml@tu\gml@yes{\bf\gml@pagenumber{#1}}\else\gml@pagenumber{#1}\fi
  \else\gml@tv\fi}
\def\gmlsee#1#2#3{\unskip.\ {\it#2}\ #3}

% ---- the document --------------------------------------------------------------------------------------
\def\gmlbegindocument{\gml@readrefs\gml@tabdefault
  \everypar{\gml@everypar}\rm\gml@apply}
\def\gmlenddocument{\par\vfill\supereject\gml@checkrefs}

% ---- tag set and local overrides -----------------------------------------------------------------------
\ifx\gmltagset\undefined\else\input\gmltagset\relax\fi
\openin\gml@in=gmllocal.tex
\ifeof\gml@in\closein\gml@in\else\closein\gml@in\input gmllocal \fi
\catcode`\@=\gmlatcode
\endinput
